\documentclass[11pt,a4paper]{article}
\usepackage[a4paper,margin=23mm,includeheadfoot,headheight=15pt,headsep=15pt,footskip=22pt]{geometry}
\usepackage[nopatch=footnote]{microtype}
\usepackage{amsmath,amssymb,amsthm,mathtools,aliascnt}
\usepackage{booktabs,longtable,array,tabularx,enumitem,fancyhdr,graphicx}
\usepackage[table]{xcolor}\usepackage{setspace,needspace,listings,caption,titlesec,etoolbox}
\PassOptionsToPackage{hyphens}{url}\usepackage{hyperref}\usepackage[nameinlink,noabbrev]{cleveref}\usepackage{bookmark}
\definecolor{accent}{HTML}{137668}\definecolor{muted}{HTML}{536568}
\frenchspacing\emergencystretch=3em\clubpenalty=10000\widowpenalty=10000\raggedbottom
\titleformat{\section}{\sffamily\bfseries\large\color{accent}}{\thesection}{.7em}{}
\titleformat{\subsection}{\sffamily\bfseries\normalsize}{\thesubsection}{.7em}{}
\titlespacing*{\section}{0pt}{1.4\baselineskip}{.55\baselineskip}
\captionsetup{font=small,labelfont={sf,bf},labelsep=period}
\pretocmd{\section}{\Needspace{7\baselineskip}}{}{}\pretocmd{\subsection}{\Needspace{5\baselineskip}}{}{}
\allowdisplaybreaks\setstretch{1.055}\setlength{\parskip}{.3em}\setlength{\parindent}{1.35em}
\setlist[itemize]{topsep=.35em,itemsep=.15em,leftmargin=1.6em}\setlist[enumerate]{topsep=.35em,itemsep=.15em,leftmargin=1.8em}
\graphicspath{{figures/}}
\hypersetup{pdftitle={Projector-Resolved Seed Calculus II: Refinement, Weighted Transfer and Integral Separation},pdfauthor={Komorebi P.},colorlinks=true,linkcolor=accent,citecolor=accent,urlcolor=accent}
\pagestyle{fancy}\fancyhf{}
\fancyhead[L]{\sffamily\fontsize{8}{10}\selectfont\color{muted}KOMOREBI MATHEMATICS / PAPER II}
\fancyhead[R]{\sffamily\fontsize{8}{10}\selectfont\color{muted}TRADITIONAL PAPER}
\fancyfoot[L]{\sffamily\fontsize{8}{10}\selectfont\color{muted}Projector-resolved seed calculus}\fancyfoot[R]{\sffamily\small\color{accent}\thepage}
\lstset{basicstyle=\ttfamily\footnotesize,breaklines=true,frame=single,showstringspaces=false,columns=fullflexible,keepspaces=true}
\newtheorem{theorem}{Theorem}[section]
\newaliascnt{proposition}{theorem}
\newtheorem{proposition}[proposition]{Proposition}
\aliascntresetthe{proposition}
\newaliascnt{lemma}{theorem}
\newtheorem{lemma}[lemma]{Lemma}
\aliascntresetthe{lemma}
\newaliascnt{corollary}{theorem}
\newtheorem{corollary}[corollary]{Corollary}
\aliascntresetthe{corollary}
\theoremstyle{definition}
\newaliascnt{definition}{theorem}
\newtheorem{definition}[definition]{Definition}
\aliascntresetthe{definition}
\newaliascnt{example}{theorem}
\newtheorem{example}[example]{Example}
\aliascntresetthe{example}
\newaliascnt{standardfact}{theorem}
\newtheorem{standardfact}[standardfact]{Standard Fact}
\aliascntresetthe{standardfact}
\theoremstyle{remark}
\newaliascnt{remark}{theorem}
\newtheorem{remark}[remark]{Remark}
\aliascntresetthe{remark}

\crefname{theorem}{theorem}{theorems}
\crefname{proposition}{proposition}{propositions}
\crefname{lemma}{lemma}{lemmas}
\crefname{corollary}{corollary}{corollaries}
\crefname{definition}{definition}{definitions}
\crefname{example}{example}{examples}
\crefname{standardfact}{standard fact}{standard facts}
\crefname{remark}{remark}{remarks}

\newcommand{\bbk}{\Bbbk}
\newcommand{\OO}{{\mathcal O}}
\newcommand{\Acal}{\mathcal A}
\newcommand{\Bcal}{\mathcal B}
\newcommand{\Pcal}{\mathcal P}
\newcommand{\Qcal}{\mathcal Q}
\newcommand{\Scal}{\mathcal S}
\newcommand{\Rcal}{\mathcal R}
\newcommand{\End}{\operatorname{End}}
\newcommand{\im}{\operatorname{im}}
\newcommand{\Ker}{\operatorname{Ker}}
\newcommand{\coker}{\operatorname{coker}}
\newcommand{\Frac}{\operatorname{Frac}}
\newcommand{\length}{\operatorname{length}}
\newcommand{\val}{\operatorname{val}}
\newcommand{\wt}{\operatorname{wt}}
\newcommand{\Supp}{\operatorname{Supp}}
\newcommand{\id}{\mathrm{id}}
\newcommand{\diag}{\operatorname{diag}}
\newcommand{\Res}{\operatorname{Res}}
\newcommand{\lcmop}{\operatorname{lcm}}
\newcommand{\expn}{\operatorname{exp}}

\providecommand{\Seed}{\mathfrak S}
\providecommand{\cX}{\mathcal X}
\providecommand{\cA}{\mathcal A}
\providecommand{\cB}{\mathcal B}
\providecommand{\cC}{\mathcal C}
\providecommand{\cD}{\mathcal D}
\providecommand{\cF}{\mathcal F}
\providecommand{\cG}{\mathcal G}
\providecommand{\cO}{\mathcal O}
\providecommand{\cP}{\mathcal P}
\providecommand{\cR}{\mathcal R}
\providecommand{\cS}{\mathcal S}
\providecommand{\cT}{\mathcal T}
\providecommand{\End}{\operatorname{End}}
\providecommand{\Hom}{\operatorname{Hom}}
\providecommand{\im}{\operatorname{im}}
\providecommand{\Ker}{\operatorname{Ker}}
\providecommand{\Spec}{\operatorname{Spec}}
\providecommand{\Res}{\operatorname{Res}}
\providecommand{\Hol}{\operatorname{Hol}}
\providecommand{\Tr}{\operatorname{Tr}}
\providecommand{\diag}{\operatorname{diag}}
\providecommand{\id}{\mathrm{id}}
\providecommand{\bbk}{\mathbb{k}}
\providecommand{\N}{\mathbb N}
\providecommand{\Z}{\mathbb Z}
\providecommand{\C}{\mathbb C}
\providecommand{\Q}{\mathbb Q}
\providecommand{\F}{\mathbb F}
\providecommand{\OO}{{\mathcal O}}
\providecommand{\Frac}{\operatorname{Frac}}
\providecommand{\val}{\operatorname{val}}
\providecommand{\wt}{\operatorname{wt}}
\providecommand{\lcm}{\operatorname{lcm}}
\providecommand{\packet}[2]{#1_{#2}}
\newcolumntype{L}[1]{>{\raggedright\arraybackslash}p{#1}}
\newcolumntype{Y}{>{\raggedright\arraybackslash}X}


\theoremstyle{definition}\newtheorem{construction}[theorem]{Construction}\newtheorem{certificate}[theorem]{Numerical Certificate}
\theoremstyle{remark}\newtheorem{boundary}[theorem]{Boundary}

\usepackage{tikz,multicol}
\usetikzlibrary{arrows.meta,decorations.pathreplacing,patterns,calc}
\usepackage{komorebi-monograph}
\renewcommand{\KomorebiKind}{Mathematical text}
\input{generated/current_edition.tex}
\KomorebiSetup
\input{atlas_macros.tex}
\usepackage{titlesec}
\input{core-layout.tex}

% Publication-design overrides. Mathematical bodies retain their source text.
\CoreLayout
\colorlet{accent}{pubTeal}\colorlet{teal}{pubTeal}\colorlet{ink}{pubInk}\colorlet{muted}{pubMuted}\colorlet{pale}{pubWash}\colorlet{light}{pubWash}
\makeatletter
\@ifundefined{KomorebiCover}{}{\renewcommand{\KomorebiCover}{%
\clearpage\hypersetup{pageanchor=false}\thispagestyle{empty}\begingroup\setlength{\parindent}{0pt}%
{\sffamily\fontsize{8.5}{11}\selectfont\color{pubTeal}REALIZATION CALCULUS\par}%
\vspace{27pt}%
{\sffamily\fontsize{9}{12}\selectfont\color{pubTeal}\MakeUppercase{\KomorebiSeries}\par}%
\vspace{22pt}%
{\normalfont\fontsize{29}{35}\selectfont\color{pubNavy}\raggedright\KomorebiTitle\par}%
\vspace{17pt}%
{\normalfont\fontsize{17}{23}\selectfont\color{pubMuted}\raggedright\KomorebiSubtitle\par}%
\vspace{22pt}\noindent\textcolor{pubCopper}{\rule{42pt}{.8pt}}\par%
\vspace{20pt}%
\@ifundefined{KomorebiCoverBlurb}{}{\normalfont\fontsize{11.4}{16.2}\selectfont\color{pubMuted}\raggedright\KomorebiCoverBlurb\par}%
\vfill%
{\normalfont\fontsize{14}{18}\selectfont\color{pubInk}Komorebi P.\par}%
\vspace{9pt}{\sffamily\fontsize{9}{13}\selectfont\color{pubMuted}\KomorebiKind\par\KomorebiEdition\par}%
\endgroup\clearpage\pagenumbering{roman}\hypersetup{pageanchor=true}\color{pubInk}\normalfont\pagestyle{empty}}}
\makeatother

\begin{document}
\KomorebiCover
\KomorebiOpeningNote{\textbf{Reading route.} Exact packet coordinates and residue selectors lead to weighted inversion and rational walk series. Lattice gluing then gives separation weights and precision delays. Smith congruences and operator defects connect these formulas to finite computation.

\medskip
\textbf{Companions.} The worked plates illustrate the constructions of this chapter. The formal companion in the back matter distinguishes the chapter-wide correspondence from the foundational 25-row walkthrough, with explicit evidence and scope.}
\begin{abstract}

A complete family of idempotents gives exact coordinates, but a single family does not record scale or integral visibility.  We study towers in which every coarse idempotent is the sum of its fine descendants.  The coordinate changes are invertible and coherent, and every operator block is the sum of its refined blocks.  Residue selectors give a model over arbitrary rings.  When roots of unity are available they agree with character averages, and coprime moduli factor through the Chinese remainder theorem.  In a complete filtered algebra, a diagonal operator plus positive weight transfers has an inverse obtained by summing admissible weighted paths.  The formula truncates exactly at finite precision, while finite transfer graphs have rational path series.  For a rational spectral decomposition over a discrete valuation ring, resultants give denominator bounds.  The exact denominator is the largest Smith exponent of the lattice gluing module.  This yields a necessary and sufficient precision delay, monotonicity under refinement, and a full congruence solving formula.  We close with exact formulas for the propagation of a coarsening defect through powers, resolvents, holomorphic functions, exponentials, and multilevel towers.
\end{abstract}
\noindent The exposition distinguishes general written proofs, formal specifications, and finite computational examples. The worked plates supply concrete models; the formal companion identifies the exact declarations and hypotheses that support each mapped statement.
\tableofcontents\listoffigures\listoftables\clearpage
\pagenumbering{arabic}
\section{Introduction}

A complete family of pairwise annihilating idempotents decomposes a module into exact coordinate pieces.  Many problems require more than one such decomposition.  A signal may be split first into two residue classes and then into four.  A filtered operator may be resolved into transfers that appear at different valuation depths.  A rational spectral splitting may exist over the fraction field while remaining invisible on the original lattice.  These are all refinement problems, but they involve different notions of scale.

The first part of the article concerns exact refinement.  If a fine packet family refines a coarse family, then each coarse idempotent is the sum of its fine descendants.  The coordinate spaces at the two levels are canonically isomorphic because both represent the whole module.  The refinement relation adds more information.  It tells us which fine coordinates belong to each coarse coordinate, and it gives a compatible decomposition of every operator block.  This is the algebra behind polyphase coordinates and nested filter banks.  Standard treatments of those subjects may be found in \cite{Vaidyanathan1993,Daubechies1992}.

The second part studies inversion in a complete filtered algebra.  Suppose an operator has the form
\[
T=D+N,
\]
where $D$ is invertible and every packet transfer in $N$ raises the filtration by a positive amount.  The ordinary Neumann series can then be expanded into weighted paths.  The weight is not the number of edges.  It is the total filtration gain along the path.  At precision $L$, every path of weight at least $L$ vanishes.  A finite transfer graph also gives a rational generating function for the number of paths of each total weight.  The relevant generating function facts are classical and are discussed in \cite{Stanley2012}.

The third part concerns integral separation.  Let $V$ be a vector space over the fraction field of a discrete valuation ring $\OO$, and let $\Lambda$ be an $\OO$ lattice in $V$.  A rational projector $E$ need not preserve $\Lambda$.  Its separation weight is the least $s$ such that
\[
\pi^sE\Lambda\subseteq\Lambda.
\]
A Bézout construction gives an upper bound through a polynomial resultant.  That bound need not be exact because it does not see the position of the lattice.  The exact answer is controlled by the finite length module
\[
G_{\Lambda,\Pcal}
=
\Lambda\Big/\bigoplus_{\alpha}(\Lambda\cap V_\alpha),
\]
where $V=\bigoplus_\alpha V_\alpha$ is the rational packet decomposition.  The largest Smith exponent of this gluing module is exactly the common denominator needed to make every packet projector integral.  This interpretation also proves that separation weight cannot decrease under refinement.  Background on local fields, Smith normal form, and matrix polynomials may be found in \cite{Serre1979,Newman1972,GohbergLancasterRodman2009}.

The last part records how a defect between fine and coarse operators propagates.  If
\[
\Delta=CT_f-T_cC,
\]
then exact formulas describe the corresponding defect for powers, resolvents, holomorphic functions, exponentials, and a chain of several refinement levels.  These formulas are elementary, but collecting them in one place makes the later analytic uses transparent.  The spectral and resolvent background follows the usual operator theory in \cite{Kato1995}.

Every theorem in this paper concerns linear operators. An exact refinement retains the full object and changes its coordinates. Questions about nonlinear operation trees, nonlinear correction, and obstruction complexes require a separate theory.

\section{Depth and compatibility}\label{sec:four_depths}

Packet, filtration, and separation depths describe three different kinds of scale. Compatibility is initially recorded by an operator defect. Keeping these notions separate makes the later formulas easier to read.

The first is packet depth.  A finer packet family has more idempotents and a more detailed address set.  The passage from residues modulo two to residues modulo four is an example.  Packet depth is combinatorial.  It says how many distinctions have been made in the coordinate system.

The second is filtration depth. A transfer in $F^r\Bcal$ vanishes modulo $F^r\Bcal$. This membership gives a lower bound on its depth. A nonzero transfer may first become visible only at a higher precision. Two paths with the same number of edges can have different filtration depths. A longer path can appear earlier when its edges have smaller weights.

The third is separation depth.  A rational packet projector may require multiplication by $\pi^s$ before it preserves the chosen lattice.  This number depends on both the rational decomposition and the lattice.  It can be zero even when the corresponding eigenvalues are close, and it can be positive even when the rational projector is simple.

The fourth notion is the compatibility defect. A fine operator and a separately chosen coarse operator need not commute with coarsening. Their discrepancy is
\[
\Delta=CT_f-T_cC.
\]
This operator records the failure of one model to represent another exactly. A chosen norm measures its size, while a chosen filtration can assign a depth. No integer compatibility depth is assigned without that additional structure.

The four notions meet in applications, but no one of them determines the others.  An exact refinement tower may have zero coordinate error while its fine spectral packets have positive lattice weight.  A filtered inverse may truncate exactly at a given depth while the coarse operator still has a nonzero compatibility defect.  The paper treats each quantity with its own theorem before combining them.

\begin{keyidea}
Packet depth changes the address of a component. Filtration depth controls when a transfer survives a precision cutoff. Separation depth clears a rational denominator on a chosen lattice. The compatibility defect records a discrepancy between two operator models. These quantities can interact, but each needs its own hypothesis and its own calculation.
\end{keyidea}

A simple two packet model already contains all four notions.  Let
\[
X=X_0\oplus X_1,
\qquad
D=\begin{pmatrix}D_0&0\\0&D_1\end{pmatrix},
\qquad
N=\begin{pmatrix}0&N_{01}\\N_{10}&0\end{pmatrix}.
\]
A finer family may split $X_0$ and $X_1$ into smaller summands.  The blocks $N_{01}$ and $N_{10}$ may enter at different filtration weights.  Over a fraction field, the rational packet spaces may fail to split the chosen lattice.  A coarse operator on one selected summand may also fail to intertwine with the full matrix.  The later sections isolate these four effects one at a time.


\input{chapters/02_coordinates.tex}
\input{chapters/03_weighted.tex}
\section{Rational spectral projectors and resultant bounds}\label{sec:resultant}

Let $\OO$ be a discrete valuation ring with fraction field $K$ and uniformizer $\pi$.  Let $\Lambda$ be a finite free $\OO$ module and put
\[
V=K\otimes_{\OO}\Lambda.
\]
Suppose $A\in\End_{\OO}(\Lambda)$.

\begin{theorem}[Bézout spectral projector]\label{thm:bezout_projector}
Assume that
\[
\chi_A(x)=f(x)g(x)
\]
with monic polynomials $f,g\in\OO[x]$ that are coprime in $K[x]$.  Choose $u,v\in\OO[x]$ and a nonzero $r\in\OO$ such that
\[
u(x)f(x)+v(x)g(x)=r.
\]
Then
\[
E_f=r^{-1}v(A)g(A)
\]
is the projector of $V$ onto $\Ker f(A)$ along $\Ker g(A)$.  In particular,
\[
E_f^2=E_f,
\qquad
rE_f\in\End_{\OO}(\Lambda).
\]
\end{theorem}

\begin{proof}
Cayley Hamilton gives
\[
f(A)g(A)=0.
\]
If $x\in\Ker f(A)$, then the Bézout identity gives
\[
v(A)g(A)x=rx,
\]
so $E_fx=x$.  If $x\in\Ker g(A)$, then $E_fx=0$.  The same identity gives the direct sum decomposition
\[
V=\Ker f(A)\oplus\Ker g(A).
\]
Indeed, every $x\in V$ is the sum of
\[
r^{-1}v(A)g(A)x\in\Ker f(A)
\]
and
\[
r^{-1}u(A)f(A)x\in\Ker g(A).
\]
Thus $E_f$ is the stated projector and is idempotent.  The formula for $rE_f$ has coefficients in $\OO$, so it preserves $\Lambda$.
\end{proof}

\begin{theorem}[Several coprime spectral packets]\label{thm:several_projectors}
Assume
\[
\chi_A(x)=\prod_{i=1}^rf_i(x),
\]
where the monic factors $f_i\in\OO[x]$ are pairwise coprime in $K[x]$.  Put
\[
g_i=\prod_{j\ne i}f_j.
\]
There are complementary projectors $E_i$ with
\[
\im E_i=\Ker f_i(A),
\qquad
\sum_iE_i=I,
\qquad
E_iE_j=0
\quad(i\ne j).
\]
One may choose $E_i$ so that
\[
\Res(f_i,g_i)E_i\in\End_{\OO}(\Lambda)
\]
up to multiplication of the resultant by a unit.
\end{theorem}

\begin{proof}
Apply \cref{thm:bezout_projector} to the factorization $\chi_A=f_ig_i$.  The resulting projector is one on $\Ker f_i(A)$ and zero on every $\Ker f_j(A)$ with $j\ne i$.  The direct sum decomposition follows from the polynomial Chinese remainder theorem.  The orthogonality and sum identities follow by evaluating the corresponding idempotents in $K[x]/(\chi_A)$.  The Sylvester Bézout relation for $(f_i,g_i)$ gives the stated denominator bound.
\end{proof}

\begin{corollary}[Resultant denominator bound]\label{cor:resultant_bound}
Under the hypotheses of \cref{thm:bezout_projector}, one may take
\[
r=\Res(f,g)
\]
up to multiplication by a unit.  Therefore
\[
\pi^{\val_\pi(\Res(f,g))}E_f\Lambda\subseteq\Lambda.
\]
\end{corollary}

\begin{proof}
The Sylvester matrix construction gives a Bézout relation with right side equal to the resultant.  Apply \cref{thm:bezout_projector}.  Standard facts about resultants may be found in \cite{Lang2002,GohbergLancasterRodman2009}.
\end{proof}

\begin{example}[A resultant bound that is not exact]\label{ex:loose_resultant}
Let
\[
A=\diag(0,\pi^s)
\]
on the standard lattice $\Lambda=\OO^2$.  The factors $x$ and $x-\pi^s$ have resultant $-\pi^s$, up to the harmless choice of resultant order, but the two spectral projectors are the integral coordinate projectors.  Their exact separation weight is zero.  The resultant sees the closeness of the two eigenvalues.  It does not see that the chosen lattice already splits along their eigenspaces.
\end{example}

\begin{interpretation}
A resultant bounds the denominator through the characteristic factors. The exact separation weight also depends on how the lattice sits inside the rational packet spaces. These are two stages of the calculation. First construct a rational projector with a known clearing denominator. Then measure the gluing of the chosen lattice to find the least denominator that actually works.
\end{interpretation}

\subsection{Evaluated certificates and explicit lattice action}

\subsubsection{Minimal evaluated B\'ezout identities}

Let $\OO$ be a DVR, $K=\Frac(\OO)$, and $A\in\End_\OO(\Lambda)$ for a finite free lattice $\Lambda$. Suppose the characteristic polynomial factors over $\OO$ as
\[
\chi_A(x)=f(x)g(x),
\]
where $f$ and $g$ are coprime over $K[x]$.

\begin{theorem}[Evaluated B\'ezout spectral projector]\label{thm:bezout-projector}
Let $F,G,U,V\in\End_K(\Lambda\otimes_\OO K)$ and let $r\in K^\times$.
Suppose
\[
UF+VG=rI,
\qquad FG=0,
\qquad FG=GF,
\qquad GU=UG.
\]
Put $N=VG$. Then
\[
\boxed{N^2=rN.}
\]
Consequently the operator
\[
\boxed{
E=r^{-1}VG
}
\]
is idempotent, acts as the identity on $\ker F$, and acts as zero on
$\ker G$. If $V$ and $G$ preserve $\Lambda$, then
\[
rE=VG\in\End_\OO(\Lambda).
\]
\end{theorem}

\begin{proof}
The B\'ezout identity gives $N=rI-UF$. The two stated commutation relations
and $FG=0$ give
\[
N(UF)=VGUF=VUGF=VUFG=0.
\]
Therefore
\[
N^2=N(rI-UF)=rN.
\]
Since $r\ne0$, $E^2=r^{-2}N^2=E$. If $Fx=0$, the evaluated B\'ezout
identity gives $VGx=rx$, hence $Ex=x$. If $Gx=0$, then $Ex=0$.
The lattice assertion follows directly from $rE=VG$.
\end{proof}

The equality $N^2=rN$ before division uses neither $r\ne0$ nor the other
pairwise commutation relations that hold in the polynomial application.
Normalizing to $E$ uses $r\ne0$. Its action on $\ker F$ uses only the B\'ezout
identity and $r\ne0$, while vanishing on $\ker G$ follows immediately from $N=VG$.
The bundled Lean certificate records all pairwise commutations as convenient
evaluated polynomial data, but the minimal theorems expose only the hypotheses
they consume.

For the polynomial application, one supplies $F=f(A)$, $G=g(A)$,
$U=u(A)$, and $V=v(A)$. A B\'ezout or resultant identity
$uf+vg=r$ evaluates to the first certificate, while
$(fg)(A)=0$ supplies the second. When $fg=\chi_A$, this last equality is
Cayley Hamilton. Standard B\'ezout, resultant, and DVR facts are recalled
from \cite{Lang2002,Serre1979}. The inherited certificate theorem begins with the evaluated
certificate. It does not infer exact action merely from nilpotence of
$f(A)$ or $g(A)$, which would be insufficient.

\begin{definition}[Separation weight]\label{def:separation-weight}
For a rational projector $E$ and a specified $\OO$ lattice $\Lambda$, provided
at least one clearing exponent exists, define
\[
s_{\pi,\Lambda}(E)
=
\min\{s\ge0:\pi^sE(\Lambda)\subseteq\Lambda\}.
\]
If the evaluated certificate arises from integral $v,g$ and
$r=u\pi^s$ with $u^{-1}$ preserving the lattice, then
\cref{thm:bezout-projector} gives the upper bound
\[
s_{\pi,\Lambda}(E)\le s=\val_\pi(r).
\]
Indeed,
\[
\boxed{\pi^sE=u^{-1}VG}
\]
preserves $\Lambda$. The minimal Lean theorem needs only the lattice action of
$u^{-1}$. The bundled wrapper for lattice units records both $u$ and $u^{-1}$.
A nonzero scalar in the fraction field need not be an integral unit.
Existence of a clearing exponent is an explicit hypothesis. Once one exists,
the least one is obtained from the least element property of the natural numbers.
Clearing exponents are upward closed provided multiplication by $\pi$
preserves the lattice.
\end{definition}

\subsubsection{Certificate-level precision delay}

\begin{theorem}[Projector precision delay]\label{thm:precision-delay}
Let $E$ be a rational projector with $\pi^sE(\Lambda)\subseteq\Lambda$. If
\[
x\equiv y\pmod{\pi^{m+s}\Lambda},
\]
then
\[
Ex\equiv Ey\pmod{\pi^m\Lambda}.
\]
Thus $m+s$ input digits suffice to determine $m$ output digits of the projected packet.
\end{theorem}

\begin{proof}
Write $x-y=\pi^{m+s}z$. Since $\pi^sE(z)\in\Lambda$,
\[
E(x-y)=\pi^m\bigl(\pi^sE(z)\bigr)\in\pi^m\Lambda.
\]
\end{proof}

The theorem gives a precision budget for delayed separation. A packet with weight $s$ can be reconstructed at the requested output precision by retaining $s$ additional input digits.


\section{Exact lattice separation}\label{sec:lattice}

The exact denominator belongs to the pair consisting of the rational packet decomposition and the integral lattice.

\Needspace{8\baselineskip}
\begin{definition}[Lattice packet data]\label{def:lattice_packet}
Let
\[
V=\bigoplus_{\alpha\in I}V_\alpha
\]
be a finite direct sum over $K$, with packet projectors $E_\alpha$.  Let $\Lambda\subseteq V$ be an $\OO$ lattice.  Put
\[
\Lambda_\alpha=\Lambda\cap V_\alpha,
\qquad
\Lambda_{\mathrm{split}}=\bigoplus_\alpha\Lambda_\alpha,
\]
and define the gluing module
\[
G_{\Lambda,\Pcal}
=
\Lambda/\Lambda_{\mathrm{split}}.
\]
The separation weight of $E_\alpha$ is
\[
s_\Lambda(E_\alpha)
=
\min\{s\geq0:\pi^sE_\alpha\Lambda\subseteq\Lambda\}.
\]
The global separation weight is
\[
s_\Lambda(\Pcal)=\max\bigl(\{0\}\cup\{s_\Lambda(E_\alpha):\alpha\in I\}\bigr).
\]
For a finitely generated torsion $\OO$ module $G$, write $\expn_\pi(G)$ for the least $s$ such that $\pi^sG=0$.
\end{definition}

The gluing module $G_{\Lambda,\Pcal}$ is a finitely generated torsion $\OO$ module, hence has finite length, because $\Lambda_{\mathrm{split}}$ and $\Lambda$ are full lattices in the same vector space. To see that the split lattice is full, clear denominators of bases of the rational summands. Each $\Lambda\cap V_\alpha$ then spans $V_\alpha$, so their direct sum spans $V$.

\begin{theorem}[Gluing module and exact separation weight]\label{thm:gluing_weight}
With the notation of \cref{def:lattice_packet}, the map
\[
\Phi:G_{\Lambda,\Pcal}
\longrightarrow
\bigoplus_{\alpha\in I}E_\alpha\Lambda/\Lambda_\alpha,
\qquad
[x]\longmapsto(E_\alpha x+\Lambda_\alpha)_\alpha
\]
is injective, and every coordinate projection of its image is surjective.  We also have
\[
\boxed{
\expn_\pi(G_{\Lambda,\Pcal})
=
s_\Lambda(\Pcal).
}
\]
If
\[
G_{\Lambda,\Pcal}
\cong
\bigoplus_{j=1}^r\OO/\pi^{a_j}\OO
\]
with $a_1\leq\cdots\leq a_r$, then
\[
s_\Lambda(\Pcal)=\max\bigl(\{0\}\cup\{a_j:1\leq j\leq r\}\bigr).
\]
\end{theorem}

\begin{proof}
If $x\in\Lambda$ maps to zero, then $E_\alpha x\in\Lambda_\alpha$ for every $\alpha$.  Since
\[
x=\sum_\alpha E_\alpha x,
\]
we have $x\in\Lambda_{\mathrm{split}}$.  Thus $\Phi$ is injective.  The coordinate map onto $E_\alpha\Lambda/\Lambda_\alpha$ is surjective by the definition of $E_\alpha\Lambda$.

Let $s=s_\Lambda(\Pcal)$.  For every $x\in\Lambda$ and every $\alpha$,
\[
\pi^sE_\alpha x\in\Lambda_\alpha.
\]
Hence
\[
\pi^sx=\sum_\alpha\pi^sE_\alpha x\in\Lambda_{\mathrm{split}},
\]
so $\pi^sG_{\Lambda,\Pcal}=0$.  This proves
\[
\expn_\pi(G_{\Lambda,\Pcal})\leq s.
\]
If $d<s$, then for some $\alpha$ there is an $x\in\Lambda$ with
\[
\pi^dE_\alpha x\notin\Lambda_\alpha.
\]
The class $\pi^d[x]$ is therefore nonzero in the gluing module.  Thus $\pi^d$ does not annihilate the gluing module, and equality follows.  The final formula is the usual description of the exponent through Smith normal form.
\end{proof}

\begin{corollary}[Total gluing and worst precision]\label{cor:length_and_weight}
If
\[
G_{\Lambda,\Pcal}
\cong
\bigoplus_{j=1}^r\OO/\pi^{a_j}\OO,
\]
then
\[
\length_{\OO}(G_{\Lambda,\Pcal})
=
\sum_{j=1}^ra_j,
\qquad
s_\Lambda(\Pcal)=\max\bigl(\{0\}\cup\{a_j:1\leq j\leq r\}\bigr).
\]
Thus the lattice index measures the total amount of gluing, while the separation weight measures the worst precision loss in one direction.
\end{corollary}

\begin{proof}
The length and exponent of a finitely generated torsion module are read directly from its Smith decomposition.  Apply \cref{thm:gluing_weight}.
\end{proof}

\begin{corollary}[Two complementary packets]\label{cor:binary_gluing}
Suppose
\[
V=V_1\oplus V_2
\]
with projector $E$ onto $V_1$ and complementary projector $I-E$.  Then
\[
\Lambda/(\Lambda_1\oplus\Lambda_2)
\cong
E\Lambda/\Lambda_1
\cong
(I-E)\Lambda/\Lambda_2.
\]
Consequently
\[
s_\Lambda(E)=s_\Lambda(I-E)
=
\expn_\pi\bigl(\Lambda/(\Lambda_1\oplus\Lambda_2)\bigr).
\]
\end{corollary}

\begin{proof}
Map $x\in\Lambda$ to $Ex+\Lambda_1$.  Its kernel consists of the $x$ for which $Ex\in\Lambda_1$.  In that case
\[
(I-E)x=x-Ex\in\Lambda\cap V_2=\Lambda_2,
\]
so the kernel is $\Lambda_1\oplus\Lambda_2$.  The map is surjective by definition.  The complementary map is handled in the same way.  The equality of weights follows from \cref{thm:gluing_weight}.
\end{proof}

\begin{corollary}[Integral splitting criterion]\label{cor:integral_split}
Every packet projector preserves $\Lambda$ if and only if
\[
\Lambda=\bigoplus_\alpha(\Lambda\cap V_\alpha).
\]
Equivalently,
\[
s_\Lambda(\Pcal)=0
\quad\Longleftrightarrow\quad
G_{\Lambda,\Pcal}=0.
\]
\end{corollary}

\begin{proof}
If every projector preserves $\Lambda$, then
\[
x=\sum_\alpha E_\alpha x
\]
expresses every $x\in\Lambda$ as a sum of vectors in $\Lambda_\alpha$.  The converse is immediate because each projector is the identity on one direct summand and zero on the others.  The final equivalence follows from \cref{thm:gluing_weight}.
\end{proof}

\subsection{Refinement and separation weight}

A finer rational packet system imposes more integral splitting conditions.

\begin{theorem}[Gluing exact sequence under refinement]\label{thm:gluing_refinement}
Let $\Qcal=(Q_\beta)$ refine $\Pcal=(P_\alpha)$ on $V$, and let $\Lambda$ be a lattice.  Define
\[
\Lambda_{\Pcal}
=
\bigoplus_\alpha(\Lambda\cap P_\alpha V),
\qquad
\Lambda_{\Qcal}
=
\bigoplus_\beta(\Lambda\cap Q_\beta V).
\]
Then
\[
\Lambda_{\Qcal}\subseteq\Lambda_{\Pcal}\subseteq\Lambda
\]
and there is a short exact sequence
\[
0
\longrightarrow
\Lambda_{\Pcal}/\Lambda_{\Qcal}
\longrightarrow
\Lambda/\Lambda_{\Qcal}
\longrightarrow
\Lambda/\Lambda_{\Pcal}
\longrightarrow0.
\]
In particular,
\[
\boxed{
s_\Lambda(\Pcal)\leq s_\Lambda(\Qcal).
}
\]
\end{theorem}

\begin{proof}
Every fine packet space lies in its parent packet space, so every element of $\Lambda_{\Qcal}$ belongs to $\Lambda_{\Pcal}$.  The displayed exact sequence is the standard quotient sequence for nested submodules.  The exponent of a quotient cannot exceed the exponent of the module from which it is obtained.  Apply \cref{thm:gluing_weight} to the two gluing modules.
\end{proof}

\begin{example}[Strict growth under refinement]\label{ex:strict_weight_growth}
Take the standard lattice $\mathbb Z_2^2$ and the involution from \cref{prop:sign_projectors}.  The one packet system $(I)$ has weight zero.  Refining it into the two rational sign packets produces weight one.  The coarse lattice is already split because there is only one summand.  The fine eigenspaces are glued together by one dyadic digit.
\end{example}

\begin{remark}[Strict growth can occur]\label{rem:strict_growth}
A coarse packet family may preserve the lattice even though a fine spectral split does not.  Refinement can reveal rational directions that are glued together integrally.  The weight records the extra depth needed to separate them.
\end{remark}

\subsection{Exact precision delay}

The separation weight has an operational meaning.  It is the exact number of extra input digits needed to recover a projected packet uniformly.

% Keep the final proof marker with its displayed conclusion.
\enlargethispage{\baselineskip}
\Needspace{8\baselineskip}
\begin{theorem}[Necessary and sufficient precision delay]\label{thm:precision_delay}
Let $E$ be a rational projector on $V$, let $\Lambda$ be a lattice, and let
\[
s=s_\Lambda(E).
\]
For every $m\geq0$, the rule
\[
[x]_{m+s}\longmapsto[Ex]_m
\]
from $\Lambda/\pi^{m+s}\Lambda$ to $V/\pi^m\Lambda$ is well defined.  Here $m$ is an integer.  If $d$ is an integer with $0\leq d<s$, then for every integer $m\geq0$ there are $x,y\in\Lambda$ such that
\[
x\equiv y\pmod{\pi^{m+d}\Lambda}
\]
but
\[
Ex\not\equiv Ey\pmod{\pi^m\Lambda}.
\]
Thus $s$ extra digits are both sufficient and necessary for a uniform precision guarantee.
\end{theorem}

\begin{proof}
If
\[
x-y=\pi^{m+s}z,
\]
then
\[
E(x-y)=\pi^m(\pi^sEz)\in\pi^m\Lambda.
\]
This proves that the rule is well defined.  Now take $d<s$.  By the minimality of $s$, there is a $z\in\Lambda$ with
\[
\pi^dEz\notin\Lambda.
\]
Set
\[
x=\pi^{m+d}z,
\qquad
y=0.
\]
Then $x\equiv y$ modulo $\pi^{m+d}\Lambda$, while
\[
E(x-y)=\pi^m(\pi^dEz)\notin\pi^m\Lambda.
\]
\end{proof}

\begin{keyidea}
The necessity statement is uniform over the lattice. A particular input may need fewer extra digits, but a smaller delay cannot recover every projected packet at the prescribed precision. This is why the least clearing exponent gives an exact precision guarantee rather than only a sufficient budget.
\end{keyidea}

\begin{example}[The same rational split with two lattices]\label{ex:lattice_dependence}
Let
\[
V=Ke_1\oplus Ke_2
\]
and let $E$ project onto $Ke_1$.  On the split lattice
\[
\Lambda_0=\OO e_1\oplus\OO e_2,
\]
the weight is zero.  Fix $s\geq1$ and consider
\[
\Lambda_s
=
\OO e_1+
\OO(\pi^{-s}e_1+e_2).
\]
Then
\[
\Lambda_s\cap Ke_1=\OO e_1,
\qquad
\Lambda_s\cap Ke_2=\pi^s\OO e_2.
\]
The gluing module is isomorphic to $\OO/\pi^s\OO$, so
\[
s_{\Lambda_s}(E)=s.
\]
The rational decomposition has not changed.  Only the lattice has changed.
\end{example}

The example can be realized by an integral operator.  Define
\[
Ae_1=0,
\qquad
Ae_2=\pi^se_2.
\]
In the basis
\[
e_1,
\quad
\pi^{-s}e_1+e_2,
\]
the matrix of $A$ is
\[
\begin{pmatrix}
0&-1\\
0&\pi^s
\end{pmatrix}.
\]
Its eigenvalue factors have resultant equal to $\pi^s$ up to a unit, and the exact projector weight on $\Lambda_s$ is also $s$.

\section{The dyadic sign split}\label{sec:dyadic}

The smallest example already shows why coefficient selectors and spectral averages must be kept distinct.

\begin{proposition}[Sign projectors over a lattice]\label{prop:sign_projectors}
Let
\[
A=
\begin{pmatrix}
1&1\\
0&-1
\end{pmatrix}
\]
act on $\mathbb Z_p^2$.  Over $\mathbb Q_p$, the sign projectors are
\[
E_+=\frac{I+A}{2}
=
\begin{pmatrix}
1&1/2\\
0&0
\end{pmatrix},
\]
\[
E_-=\frac{I-A}{2}
=
\begin{pmatrix}
0&-1/2\\
0&1
\end{pmatrix}.
\]
For every odd prime $p$, both projectors preserve $\mathbb Z_p^2$.  At $p=2$,
\[
s_2(E_+)=s_2(E_-)=1.
\]
The factors $x-1$ and $x+1$ have resultant $2$.
\end{proposition}

\begin{proof}
The displayed matrices are complementary idempotents.  The scalar $2$ is a unit in $\mathbb Z_p$ for odd $p$, so both matrices are integral there.  Over $\mathbb Z_2$, multiplying by $2$ clears the denominators, while neither projector preserves the standard lattice.  Thus both weights are one.  The resultant of the two linear factors is their root difference, which is $2$.
\end{proof}

Modulo $2$, the two eigenvalues become equal.  The rational sign packets remain distinct over $\mathbb Q_2$, but the standard lattice does not split into their intersections until one more dyadic digit is retained.

\Needspace{4\baselineskip}
\begin{proposition}[Integral character averages away from the group order]\label{prop:character_integrality}
Let $\OO$ be a discrete valuation ring, let $N$ be a unit in $\OO$, and suppose $\OO$ contains a primitive $N$th root of unity $\omega$.  If an integral operator $U$ satisfies $U^N=I$, then
\[
E_j=\frac1N\sum_{k=0}^{N-1}\omega^{-jk}U^k
\]
preserves every $U$ stable $\OO$ lattice.  At primes dividing $N$, this formula alone gives no integral projector.
\end{proposition}

\begin{proof}
Every scalar coefficient in the displayed sum belongs to $\OO$ because $N$ is a unit and $\omega$ is integral.  Each power of $U$ preserves the lattice, so their $\OO$ linear combination does as well.  When $N$ is not a unit, division by $N$ is not an operation in $\OO$, so the same argument is unavailable.
\end{proof}

At a prime dividing $N$, the factor $1/N$ no longer guarantees that the character average preserves the lattice. Cancellation may still make a particular projector integral. Coefficient residue selectors remain exact because they never divide by the group order.

\subsection{The concrete rational sign matrix}

Consider the integral involution
\[
A=
\begin{pmatrix}
1&1\\
0&-1
\end{pmatrix},
\qquad
A^2=I.
\]
Over a field of characteristic different from two, its sign projectors are
\[
E_+=\frac{I+A}{2}
=
\begin{pmatrix}
1&1/2\\
0&0
\end{pmatrix},
\]
\[
E_-=\frac{I-A}{2}
=
\begin{pmatrix}
0&-1/2\\
0&1
\end{pmatrix}.
\]
They are orthogonal idempotents and sum to $I$.

\begin{figure}[htbp]
\centering
\includegraphics[width=0.96\textwidth]{07_dyadic_delayed_separation.pdf}
\caption{The rational sign projectors for an integral involution. Their denominator has $2$-adic weight one.}
\label{fig:dyadic-projectors}
\end{figure}

\begin{proposition}[Odd-adic versus dyadic behavior]\label{prop:dyadic-sign}
For every prime $p\ne2$, every entry of $E_+$ and $E_-$ is $p$-integral,
and each projector has exact local denominator weight zero. At $p=2$, each
projector has exact local denominator weight one. Multiplication by $2$ makes
it integral, while the unscaled matrix is not integral:
\[
s_{2,\Z_2^2}(E_+)=s_{2,\Z_2^2}(E_-)=1.
\]
The characteristic factors $x-1$ and $x+1$ have resultant $2$. Its $2$-adic
valuation is one, exactly the two local projector weights above, and its
$p$-adic valuation is zero at every prime $p\ne2$.
\end{proposition}

\begin{proof}
The reduced denominators of the displayed entries are $1$ or $2$, which proves
the local integrality and minimality statements directly. The resultant is
\[
\Res(x-1,x+1)=2.
\]
Thus $v_2(2)=1$ and $v_p(2)=0$ for primes $p\ne2$.
\end{proof}

The two characteristic factors $x-1$ and $x+1$ have the same reduction in
$\F_2[x]$. The rational matrices $E_+$ and $E_-$ themselves cannot be reduced
modulo $2$, because their entries contain $1/2$. Consequently the integral
special fibre does not inherit their separate idempotent splitting, even
though the two projectors remain distinct over $\Q_2$. This is the simplest
model of delayed dyadic separation.

\begin{figure}[htbp]
\centering
\includegraphics[width=0.96\textwidth]{08_odd_prime_separation.pdf}
\caption{The denominator weight of the sign projector is nonzero only at the prime two.}
\label{fig:odd-dyadic}
\end{figure}

\Needspace{7\baselineskip}
\subsubsection{Why coefficient packets remain available}

The failure of the character formula
\[
\frac12(I\pm A)
\]
to be integral does not invalidate coefficient selectors such as even and odd residue extraction. It shows that two distinct constructions must be distinguished:
\[
\boxed{
\text{integral address selector}
\quad\neq\quad
\text{character-average spectral projector}
}
\]
when the group order is not a unit. The first can remain exact over $\Z_2$. The second can acquire a denominator weight.


\section{Smith normal form and congruence lifting}\label{sec:smith}

The gluing module gives the exact common denominator of packet projectors.  Smith normal form gives the corresponding coordinate by coordinate solvability and branching rule.

\begin{standardfact}[Smith normal form over a discrete valuation ring]\label{fact:smith}
Let $M\in M_{m,n}(\OO)$.  There are invertible matrices
\[
U\in\operatorname{GL}_m(\OO),
\qquad
V\in\operatorname{GL}_n(\OO)
\]
such that
\[
UMV
=
\diag(\pi^{s_1},\ldots,\pi^{s_r},0,\ldots,0),
\]
where
\[
0\leq s_1\leq\cdots\leq s_r
\]
and $r$ is the rank over $K$.  The exponents are uniquely determined.  See \cite{Newman1972}.
\end{standardfact}

\begin{proposition}[Computing a packet weight from an inclusion matrix]\label{prop:weight_from_inclusion}
Choose an $\OO$ basis of $\Lambda$ and an $\OO$ basis of $\Lambda_{\mathrm{split}}$.  Let $B$ be the matrix of the inclusion
\[
\Lambda_{\mathrm{split}}\hookrightarrow\Lambda.
\]
If the Smith form of $B$ is
\[
\diag(\pi^{a_1},\ldots,\pi^{a_n}),
\qquad
0\leq a_1\leq\cdots\leq a_n,
\]
then
\[
G_{\Lambda,\Pcal}
\cong
\bigoplus_{j=1}^n\OO/\pi^{a_j}\OO,
\qquad
s_\Lambda(\Pcal)=\max\bigl(\{0\}\cup\{a_j:1\leq j\leq n\}\bigr).
\]
\end{proposition}

\begin{proof}
The gluing module is the cokernel of the inclusion matrix $B$.  Smith normal form gives its invariant factors.  The weight formula follows from \cref{thm:gluing_weight}.
\end{proof}

\begin{theorem}[Smith weighted congruence solve]\label{thm:smith_congruence}
Let $M\in M_{m,n}(\OO)$ have Smith exponents $s_1,\ldots,s_r$.  Assume the residue field is finite of cardinality $q$, fix $L\geq1$, and consider
\[
Mx\equiv b\pmod{\pi^L}.
\]
Put
\[
\widetilde b=Ub
\]
for a Smith matrix $U$ as in \cref{fact:smith}.  The congruence is solvable if and only if
\[
\widetilde b_i\equiv0
\pmod{\pi^{\min(s_i,L)}}
\qquad(1\leq i\leq r)
\]
and
\[
\widetilde b_i\equiv0
\pmod{\pi^L}
\qquad(r<i\leq m).
\]
When it is solvable, the number of solutions in $(\OO/\pi^L\OO)^n$ is
\[
\boxed{
q^{(n-r)L+\sum_{i=1}^r\min(s_i,L)}.
}
\]
\end{theorem}

\begin{proof}
Set $x=Vy$.  Multiplying the congruence by $U$ reduces it to independent equations
\[
\pi^{s_i}y_i\equiv\widetilde b_i\pmod{\pi^L}
\qquad(1\leq i\leq r)
\]
and zero equations in the remaining rows.  Multiplication by $\pi^{s_i}$ on $\OO/\pi^L\OO$ has image equal to the elements divisible by $\pi^{\min(s_i,L)}$ and kernel of size
\[
q^{\min(s_i,L)}.
\]
Each of the $n-r$ unused variables is free and contributes $q^L$ choices.  Multiplying the independent counts gives the formula.
\end{proof}

\begin{corollary}[Obstruction and branching data]\label{cor:smith_obstruction}
The transformed entries of $\widetilde b$ that fail the stated divisibility conditions are the complete finite congruence obstructions.  When all obstructions vanish, the Smith exponents determine the number of solution branches at depth $L$.
\end{corollary}

\begin{proof}
This is exactly the solvability and counting statement of \cref{thm:smith_congruence} written in invariant coordinates.
\end{proof}

\begin{example}[Three weighted directions]\label{ex:smith_three}
Take $p=2$, $L=6$, and the diagonal weights
\[
(0,1,3).
\]
The right side
\[
(1,2,8)
\]
is admissible.  The three coordinate kernels have sizes
\[
1,
\quad2,
\quad8,
\]
so there are $16$ solutions.  Replacing the final entry $8$ by $4$ violates the required divisibility by $2^3$ and makes the system unsolvable.
\end{example}

\subsection{Concrete diagonal fibres over integer quotients}

Let $M$ be a matrix over a DVR. Smith normal form gives invertible matrices $U,V$ such that
\[
UMV
=
\diag(\pi^{s_1}u_1,\ldots,\pi^{s_r}u_r,0,\ldots,0),
\]
where each $u_i$ is a unit and
\[
0\le s_1\le\cdots\le s_r.
\]
The exponents $s_i$ are directional lifting weights.

\begin{theorem}[Concrete Smith-weighted congruence solve]\label{thm:smith-solve}
Let $p>1$ and consider the diagonal system over
$(\Z/(p^L)\Z)^d$
\[
\diag(p^{s_1},\ldots,p^{s_d})x=b
\pmod{p^L}.
\]
The system is solvable if and only if
\[
p^{\min(s_i,L)}\mid (b_i)_{\mathrm{val}}
\]
for every $i$, where $(b_i)_{\mathrm{val}}$ is the canonical natural
representative. Conditional on this admissibility condition, the number of
solutions is
\[
\boxed{
p^{\sum_{i=1}^d\min(s_i,L)}.
}
\]
If any coordinate is inadmissible, the solution set is empty and its
cardinality is zero.
\end{theorem}

\begin{proof}
Lean proves the one-coordinate arithmetic directly over
$\mathbb Z/p^L\mathbb Z$ for every $p>1$:
\[
\bigl(\exists x\in\mathbb Z/p^L\mathbb Z:\quad p^s x=b\pmod{p^L}\bigr)
\quad\Longleftrightarrow\quad
p^{\min(s,L)}\mid b_{\mathrm{val}},
\]
where $b_{\mathrm{val}}$ is the canonical natural representative, and every
admissible fibre has exactly $p^{\min(s,L)}$ elements. The coordinates are
independent, so the product equivalence gives the displayed global count. If
one coordinate fibre is empty, the global fibre is empty. Lean proves both the
global solvability equivalence and the conditional count, including
cardinality zero on the inadmissible branch.
\end{proof}

For an arbitrary DVR $\OO$ with residue field cardinality $q$, the analogous
formula
\[
q^{\sum_i\min(s_i,L)}
\]
follows from the solvability and fibre count for each coordinate. The formal
correspondence records direct proofs of this DVR coordinate arithmetic and
the rectangular transport as separate targets.
The inherited abstract product interface remains available when the coordinate
facts are supplied as certificates. A Smith form and its changes of basis are
inputs to the transport statement; it does not provide an extracted
diagonalization routine. See \cite{Serre1979,Lang2002}.

If a Smith form is rank deficient, a literal zero diagonal factor requires
the corresponding target coordinate to vanish modulo $p^L$ and then
contributes all $p^L$ possible coordinate values. In the concrete quotient it
can equivalently be represented by any weight $s_i\ge L$, since
$p^{s_i}=0$ in $\Z/(p^L)\Z$. The theorem then gives exactly that condition and
factor. The numbered theorem takes the diagonal data as its input.

\begin{example}
At $p=2$, depth $L=6$, and weights $(0,1,3)$, the right side $(1,2,8)$ is admissible. The coordinate kernels have sizes
\[
1,2,8,
\]
so the system has $16$ solutions. Replacing $8$ by $4$ in the final coordinate violates the divisibility condition and makes the system unsolvable.
\end{example}

\begin{interpretation}
The gluing calculation takes the largest Smith weight because every packet must be separated to the requested precision. The congruence calculation retains the individual directional weights. A coordinate with weight $s_i$ contributes a solvability threshold and $\min(s_i,L)$ to the branching exponent. The independent coordinate counts multiply, so these truncated weights add in the exponent.
\end{interpretation}



\input{chapters/05_defects.tex}
\section{Worked realizations in the common operator language}

\subsection{Signal polyphase coordinates}

Let $x=(x_n)_{n\in\Z}\in c_{00}(\Z;R)$ and let the seed be $\delta_0$
under the bilateral shifts $(S^r\delta_0)_n=\delta_{n,r}$. This makes the
algebraic span and every displayed sum finite. The $N$ packets are
\[
x_j[m]=x[Nm+j].
\]
Synthesis interleaves these sequences. A linear filter with a finite impulse
response that is invariant under shifts becomes an $N\times N$ matrix of
transfer blocks. Refining from $N$ to $M$ with $N\mid M$ splits each coarse
phase into its descendants without changing the reconstructed signal. The
same formulas extend to standard sequence completions when the synthesis
series converges. This is the usual polyphase viewpoint from multirate filter
banks \cite{Vaidyanathan1993}.

\subsection{Sparsity of transfer blocks}

A banded operator need not remain banded after residue reordering, but its
nonzero couplings become a structured packet graph. With the convention
$(S_rx)_n=x_{n-r}$, the only active transfer is
\[
P_{N,j+r}S_rP_{N,j},
\]
with indices taken modulo $N$. More general finite stencils produce a finite set of active phase differences.

\begin{figure}[htbp]
\centering
\includegraphics[width=0.96\textwidth]{10_transfer_block_sparsity.pdf}
\caption{Norms of packet blocks for a finite stencil. The original operator is exactly reconstructed from a sparse transfer pattern.}
\label{fig:block-sparsity}
\end{figure}

\subsection{A short realization in Smith chart coordinates}\label{sec:smith-realization}

\begin{interpretation}
The Smith chart used here is a coordinate diagram for impedance and reflection. It is different from the Smith normal form used earlier to study integral congruences. The connection within this paper is the way a change of coordinates makes an action easier to describe. Here a rotation in reflection coordinates becomes a fractional linear map in impedance coordinates. The sampled phase orbit is a geometric example and does not supply another linear projector decomposition.
\end{interpretation}

Let $Z$ and $Z_0\ne0$ be complex impedances, put $z=Z/Z_0$, and assume $z\ne-1$. Write
\[
\Gamma=\frac{z-1}{z+1}.
\]
On the open reflection disk $|\Gamma|<1$, multiplication by any phase $\omega_N^k$ stays in the disk and has a finite inverse-chart image. The phase action is
\[
\Gamma\mapsto\omega_N^k\Gamma.
\]
Outside the open disk, the same chart formula requires $1-\omega_N^k\Gamma\ne0$. Pulling the phase action back on this domain gives
\[
T_{N,k}(z)
=
\frac{(1+\omega_N^k)z+(1-\omega_N^k)}
{(1-\omega_N^k)z+(1+\omega_N^k)}.
\]
For $N=2$ and $k=1$,
\phantomsection
\[
T_{2,1}(z)=\frac1z,
\qquad z\ne0,-1.
\label{eq:smith-inversion}
\]
Thus the nontrivial action at the first phase level exchanges normalized
impedance and admittance. Calling these values ``packets'' is geometric
terminology for the sampled orbit. It does not assert that the two values
form linear idempotent ranges. Higher refinements at powers of $2$ sample the
same reflection circle at progressively finer electrical phases. Compare the
standard Smith chart conventions in \cite{pozar2012}.

\begin{figure}[htbp]
\centering
\includegraphics[width=0.96\textwidth]{11_smith_polyphase_orbit.pdf}
\caption{Nested phase orbits at powers of $2$ on the Smith reflection disk. The full treatment of the Smith chart is deferred to the signal paper.}
\label{fig:smith-orbit}
\end{figure}

The symmetric weight
\[
w_{\mathrm{Smith}}(\Gamma)
=
\frac1{|1-\Gamma|^2}
+
\frac1{|1+\Gamma|^2}
\]
indicates where small errors in reflection coordinates produce large impedance or admittance errors.

\begin{figure}[htbp]
\centering
\includegraphics[width=0.96\textwidth]{12_smith_adaptive_weight.pdf}
\caption{A symmetric refinement weight concentrated near the boundary points for open and short circuits.}
\label{fig:smith-weight}
\end{figure}


\section{Finite computational certificates}
The following results and tables record finite checks of the implementation. They do not certify the general lattice or analytic statements in Lean. The accompanying verification files record the execution evidence and tested ranges.


All numerical work in this paper is finite and reproducible. It tests the
Python implementation of identities proved above or checked against
certificates. It does not replace those arguments or the Lean kernel checks.
The recorded implementation run covers 48 tests, including contracts for invalid input, refinement
transitivity, aggregation of transfer blocks, Fourier/DFT agreement, Smith
fibre counts by exhaustive enumeration, and consistency of the stored JSON
and CSV records.

\begin{table}[htbp]
\centering
\caption{Representative numerical and symbolic checks.}
\begin{tabularx}{\textwidth}{L{0.54\textwidth}Y}
\toprule
Check & Result \\
\midrule
Packet synthesis residuals for moduli $2,3,4,6,8,12$ & $0$ \\
Nested refinement residuals for $2\to4$, $2\to8$, $3\to6$, $3\to12$, $4\to8$, $4\to12$, $6\to12$ & $0$ \\
CRT residuals for dyadic and odd moduli for $4\times3$, $8\times3$, $4\times5$, $8\times5$ & $0$ \\
Maximum finite Fourier projector residual & $6.60\times10^{-15}$ \\
Weighted inverse residual modulo $2^{10}$ & $0$ \\
Weighted inverse residual modulo $3^8$ & $0$ \\
Weighted inverse residual modulo $5^7$ & $0$ \\
Difference between filtered and direct modular inverses & $0$ \\
Transfer zeta coefficients versus direct walk enumeration through weight $30$ & exact match \\
Dyadic sign projector resultant weight & $v_2(2)=1$ \\
Sign projector resultant weight at the odd prime $3$ & $v_3(2)=0$ \\
Power defect telescoping identity residual & at most $6.94\times10^{-18}$ \\
Resolvent defect identity residual & $1.09\times10^{-17}$ \\
Smith $N=2$ phase versus $z\mapsto1/z$ & residual below $1.4\times10^{-16}$ \\
\bottomrule
\end{tabularx}
\end{table}

\begin{certificate}[Weighted inverse]
For the tested diagonal-plus-transfer matrices, the filtered inverse agrees exactly with direct modular inversion at all tested precisions:
\[
2^{10},\qquad3^8,\qquad5^7.
\]
The maximum modular residual is zero in every case.
\end{certificate}

\begin{certificate}[Transfer zeta]
The rational function
\[
\frac{-t^3}{t^9-t^7-t^6-t^5+2t^4+t^3+t^2-1}
\]
was expanded through degree $30$. Every coefficient agrees with direct weighted-walk enumeration.
\end{certificate}

\begin{certificate}[Dyadic separation]
The code verifies exactly that the two rational sign projectors are idempotent, orthogonal, and complementary, that scaling by $2$ makes both integral, and that the same projectors are already integral modulo every tested odd prime.
\end{certificate}


\section{Scope and limits}\label{sec:limits}

The exact statements above have several deliberate boundaries.

First, refinement changes coordinates while preserving reconstruction. If a numerical method discards fine packets, the resulting error must be estimated separately.

Second, the weighted inverse is linear.  A nonlinear map creates higher order interactions among packet coordinates.  Those interactions are not contained in the path formula of \cref{thm:weighted_inverse}.

Third, a resultant is only a denominator bound.  The exact weight depends on the lattice and is given by \cref{thm:gluing_weight}.  The two examples in \cref{ex:loose_resultant,ex:lattice_dependence} show that the same eigenvalue separation can give different integral packet weights.

Fourth, the dyadic sign example is a model rather than a classification of all phenomena at the prime two.  Character averages, trace splittings, quadratic refinements, and deformation projectors must each be checked on their own lattices.

Finally, the defect formulas do not prove that a chosen coarse operator is a good approximation.  They say exactly how a known one step defect propagates.  A useful approximation theorem must also bound that defect in the problem under study.



\appendix
\section{Notation table}

\begin{longtable}{L{0.28\textwidth}L{0.64\textwidth}}
\toprule
Symbol & Meaning \\
\midrule
\endhead
$P_{\ell,\alpha}$ & packet projector at refinement level $\ell$ \\
$I_\ell$ & packet address set at level $\ell$ \\
$\beta\succ\alpha$ & fine packet $\beta$ is a child of coarse packet $\alpha$ \\
$\cA_\ell,\cS_\ell$ & packet analysis and synthesis maps \\
$J_{\ell+1,\ell}$ & exact coordinate refinement map \\
$C_{\ell,\ell+1}$ & exact coordinate regrouping from fine to coarse \\
$T^{(\ell)}_{\beta\alpha}$ & transfer block $P_{\ell,\beta}TP_{\ell,\alpha}$ \\
$P_{N,j}$ & residue or phase selector for address $j\pmod N$ \\
$w(e)$ & positive filtration weight of a transfer edge \\
$W(\gamma)$ & cumulative weight of a transfer word or graph walk \\
$A(t)$ & polynomial weighted adjacency matrix \\
$Z_{s,t}(t)$ & weighted transfer-zeta function from vertex $s$ to vertex $t$ \\
$s_{\pi,\Lambda}(E)$ & separation weight of a rational projector on the chosen lattice, when a clearing exponent exists \\
$\Delta$ & linear refinement defect for one step $CT_f-T_cC$ \\
\bottomrule
\end{longtable}

\section{Algorithmic forms}\label{sec:algorithms}

\subsection{Residue analysis and synthesis}\label{alg:residue}

\begin{lstlisting}
def analyze_residues(x, N):
    require N > 0
    return [x[j::N] for j in range(N)]

def synthesize_residues(packets):
    N = len(packets)
    require packet lengths fit one N-way interleaving
    y = empty(sum(len(p) for p in packets))
    for j, packet in enumerate(packets):
        y[j::N] = packet
    return y
\end{lstlisting}

The verified contract is
$\operatorname{synthesize}(\operatorname{analyze}(x,N))=x$. The Lean
counterpart uses packet coordinates restricted to projector ranges and proves
both inverse laws. The finite Python routine checks the concrete array
realization.

\subsection{Weighted inverse at finite precision}\label{alg:weighted-inverse}

\begin{lstlisting}
Dinv = inverse(D)                       # certified on both sides
edge[e] = Dinv @ H_edge[e]
require Dinv @ H == sum(edge[e] for e in edges)
require edge[e] lies in filtration[weight[e]]
m = certified_lower_bound(weight)       # 1 <= m <= weight[e]
N = (L + m - 1) // m                    # so L <= m*N
S = 0
for n in range(N):
    for word in edge_tuples(edges, n):
        if sum(weight[e] for e in word) < L:
            S += (-1)**n * ordered_product(edge[e] for e in word)
return q_L(S @ Dinv)
\end{lstlisting}

Lean proves that the mapped pruned sum is exactly the mapped finite geometric
inverse. The lower bound $m$ gives $m|\gamma|\le W(\gamma)$, so
$L\le mN$ derives $q_L(A^N)=0$ from the filtration certificate and makes the
returned value an inverse on both sides of $q_L(D+H)$. The Lean lemma
$L\le m\lceil L/m\rceil$ justifies the displayed executable cutoff for
$m>0$. Taking $m=1$ and $N=L$ gives
the unconditional cutoff for positive weights. The Python implementation checks the
DVR matrix specialization against direct modular inversion.

\subsection{Walk construction for transfer zeta}\label{alg:path-zeta}

\begin{lstlisting}
def coefficient(G, start, target, n):
    total = 0
    for length in range(n + 1):
        for word in edge_tuples(G, length):
            if is_walk(word, start, target) and weight(word) == n:
                total += 1
    return total

A[i,j] = sum(t**weight(e) for distinct edges e from i to j)
T = I - A
Q = det(T)
H = adjugate(T)
P = H[start, target]
assert Q.coeff(0) == 1
Z = coefficient_series(G, start, target)
prove Q * Z == P
\end{lstlisting}

Positive weights make the coefficient loop finite because every contributing
word has length at most $n$. Lean proves the enumeration correct, derives its
recurrence from the first edge, proves $(I-A)Z=I$, constructs the
determinant/adjugate inverse, and identifies the two matrices by inverse
uniqueness. Thus the normalized $QZ=P$ identity and its induced recurrence
follow from finite graph data. The Python companion repeats the exact symbolic
construction as an executable check. The Lean proof does not use it as an
oracle.

\section{Proof dependency map}

\begin{figure}[htbp]
\centering
\includegraphics[width=0.96\textwidth]{13_theorem_dependency.pdf}
\caption{The dependency structure of Chapter II. Coordinate coherence, residue/CRT, weighted inversion, automatic finite-graph transfer rationality, and integral separation are distinct branches with explicit inputs.}
\label{fig:dependency}
\end{figure}

The proof order is:
\[
\begin{aligned}
&\text{complete packet axioms}
  \Rightarrow\text{finite-packet uniqueness and coordinate coherence}\\
&\qquad\Rightarrow\text{coordinate naturality and refined transfer blocks},\\
&\text{residue selectors and a finite Fourier table}
  \Rightarrow\text{Fourier formulas},\\
&\text{coprime moduli}\Rightarrow\text{dyadic--odd formulas},\\
&\text{quotient, filtration, and edge-weight certificates}
  \Rightarrow\text{explicit pruned inverse},\\
&\text{finite positive-weight graph}
  \Rightarrow\text{semantic enumeration and first-edge recurrence}\\
&\qquad\Rightarrow\det/\operatorname{adj}\text{ rationality and recurrence},\\
&\text{minimal evaluated B\'ezout data and inverse-unit lattice action}\\
&\qquad\Rightarrow\text{separation bound, upward closure, and precision delay},\\
&\text{concrete }\mathbb Z/p^L\mathbb Z\text{ coordinate fibres}
  \Rightarrow\text{finite Smith product count},\\
&\text{one-step typed defect}
  \Rightarrow\text{powers, polynomial actions, and witnessed resolvents}.
\end{aligned}
\]

The walk branch for transfer zeta is logically independent of operator
inversion. Its rationality conclusion follows from finite graph data without
a supplied $QZ=P$ premise. Likewise, CRT factorization follows directly from
residue selectors and coprimality. It does not require the abstract refinement
theorem, although it supplies a concrete refinement family. General Smith
diagonalization and the construction of polynomial B\'ezout/resultant data
remain inputs to the formalized branches that use those certificates.


\input{chapters/07_scope.tex}
\clearpage
\section{Visual atlas}\label{app:visual-atlas}
Twenty-four worked plates follow the main constructions through explicit vectors, matrices, weighted walks, and congruence fibres. Their labels Plate II.1 to Plate II.24 stay fixed across the integrated atlas and the standalone guide. Each plate names the corresponding statement and its printed page in this monograph.

\begin{readerroadmap}[title={Three routes through the atlas}]
\textbf{Exact coordinates.} Read Plate II.1 to Plate II.8, then the finite filter in Plate II.20.

\textbf{Weights and arithmetic.} Read Plate II.9 to Plate II.19. The first group follows weighted operator words and graph walks. The second follows denominator clearing and Smith congruences.

\textbf{Realizations and defects.} Read Plate II.20 to Plate II.24. These plates separate the Smith impedance chart from Smith normal form, trace error propagation, and connect the examples to the formal contract.
\end{readerroadmap}
The matrices and finite ranges shown here are explicit models. The general proofs remain in the preceding sections. A picture or finite computation does not extend the hypotheses of a theorem or the coverage of a Lean declaration.
\clearpage
\begingroup
\setstretch{1.01}\fontsize{10}{12}\selectfont\setlength{\parskip}{3pt}\setlength{\abovedisplayskip}{6pt}\setlength{\belowdisplayskip}{6pt}
% Worked exposition inherited from the earlier edition; current navigation and scope are synchronized.

\vpage{1}{Three kinds of resolution}{CHAPTER II / A WORKED VISUAL GUIDE}{Chapter II, mathematical text: Definition 3.4, Theorem 5.10, Theorem 7.7}
Chapter II asks three different questions that can all sound like ``increase the resolution.'' This guide separates them, then follows small objects far enough that the distinctions can be computed rather than merely named.

\vfig[.96]{three_resolutions}
\platecaption{Three different operations. On the left, all coefficients survive a regrouping. In the middle, an integer weight controls which operator words can matter modulo a chosen power. On the right, a rational projector needs a denominator cleared before integral precision can be recovered.}

\head{Address resolution}
A coefficient such as $x_6$ can sit in the even packet, then the residue-$2$ packet modulo $4$, then the residue-$6$ packet modulo $8$. Its value does not become more accurate. Its \emph{address becomes more specific}, while the full object remains exactly reconstructible.

\head{Arithmetic precision}
An operator word with total filtration weight at least $L$ vanishes in the specified depth-$L$ quotient. Increasing $L$ can make additional words relevant. This is a change in the precision of the ring in which the answer is required, not a lossless rearrangement of one fixed list.

\head{Separation precision}
For the integral involution used later, $E_+=(I+A)/2$ is an exact rational projector but not an integral one at $2$. One extra input digit suffices to recover the projected difference to the requested precision. The packet is present over the fraction field even when its integral separation is delayed.

\takeaway{Keep three quantities distinct: the packet address, the precision cutoff, and the denominator weight. Finer addresses do not automatically improve arithmetic precision or remove a spectral denominator.}

\vpage{2}{An atlas of the refinement layer}{ORIENTATION / SYMBOLS WITH JOBS}{Chapter II, mathematical text: Definition 3.2, Definition 3.4, Definition 7.6}
The same letter can have different roles in different realizations. The table fixes the roles used here. Every small matrix in the guide is explicitly supplied. No diagonalization or classification theorem is inferred from a picture.

{\small
\begin{tabularx}{\linewidth}{@{}p{.15\linewidth}Yp{.28\linewidth}@{}}\toprule
\textbf{Object} & \textbf{Role} & \textbf{Worked instance}\\\midrule
$P_{N,j}$ & Retain coefficients with address $j$ modulo $N$. & Even/odd masks, then mod-$4$ and mod-$8$ masks.\\
$\mathcal A_N,\mathcal S_N$ & Split into range-restricted packet vectors, then sum them. & A permutation/interleaving of eight coordinates.\\
$J,C$ & Exact refinement and regrouping. In the defect chapter, $C$ may instead be any supplied coarsening map. & A permutation in Plate II.5 and a shear in Plate II.22.\\
$T_{ba}$ & Input packet $a$, output packet $b$: $P_bTP_a$. & Blocks of a finite cyclic filter.\\
$D,H,A$ & $T=D+H=D(I+A)$, with $A=D^{-1}H$. & Explicit integral $2\times2$ matrices.\\
$w(e),W(\gamma)$ & Positive certified edge gain and its additive word weight. & $w(a)=1$, $w(b)=2$.\\
$L$ & Precision depth in a quotient such as $\ZZ/2^L\ZZ$. & $L=5$ gives arithmetic modulo $32$.\\
$A(t),Z(t)$ & Polynomial graph adjacency and unsigned walk-count series. & Two parallel edges remain distinct.\\
$E,\Lambda,s$ & A rational projector, a specified lattice, and a clearing exponent. & $E_+$ has local weight $1$ at $2$.\\
$(s_i)$ & Directional weights of a supplied Smith diagonal. & $(1,2)$ modulo $8$ or $(0,1,3)$ across depths.\\
$\Delta$ & The failure of an independently chosen operator to intertwine: $CT_f-T_cC$. & One nonzero entry $1/16$.\\\bottomrule
\end{tabularx}}

\head{Suggested reading routes}
For exact coordinates, read \pagerefV{3} through \pagerefV{8}, then the finite filter on \pagerefV{20}. For weighted inversion and counting, read \pagerefV{9} through \pagerefV{13}. For arithmetic separation and congruence fibers, read \pagerefV{14} through \pagerefV{19}. The final pages distinguish nonlinear chart geometry, error propagation, and the formal contract.

\takeaway{Operator words carry signed matrix values. Graph walks carry unsigned counts. Smith normal form concerns linear congruences. The Smith impedance chart on Plate II.21 is a different, nonlinear construction.}

\vpage{3}{Follow one coefficient through the tower}{MODEL A / LOSSLESS ADDRESS REFINEMENT}{Chapter II, mathematical text: Theorem 3.3, Theorem 4.6}
Take a vector $x=(x_0,\ldots,x_{15})$. A packet is not a single summary number: it retains every coefficient at its selected addresses. For example,
\[
P_{2,0}x=(x_0,0,x_2,0,\ldots,x_{14},0).
\]
The rail below lists the addresses in each packet. Its display order groups descendants under their parents. The matrix example on Plate II.5 instead uses increasing residue order.

\vfig[.96]{coefficient_tower}
\platecaption{Every row contains all sixteen addresses exactly once. A heavier boundary separates packet ranges. The numbers are coefficient addresses, not new coefficient values. There is no interpolation, averaging, or omitted high-frequency tail.}

The address $6$ is even, is $2$ modulo $4$, and is $6$ modulo $8$. Its coefficient therefore lies in the successive ranges
\[
P_{8,6}X\subseteq P_{4,2}X\subseteq P_{2,0}X.
\]
The corresponding identities hold as operators:
\[
P_{2,0}=P_{4,0}+P_{4,2},\qquad
P_{4,2}=P_{8,2}+P_{8,6}.
\]
Applying either composite parent selector to $P_{8,6}x$ leaves it unchanged. Applying the odd selector kills it. This is absorption between a parent and a child, rather than an approximate resemblance between nearby scales.

\head{Why the reconstruction is unique}
If a vector is a sum of components in disjoint packet ranges, apply the projector for one range. All other components vanish and the selected component remains. Thus every child component is recovered uniquely from the sum.

\takeaway{Refinement is a sharper address book for the same coefficients. The total number of free coordinates does not change.}

\vpage{4}{Regrouping is not averaging}{MODEL A / THE MEANING OF COARSENING}{Chapter II, mathematical text: Definition 3.4, Theorem 3.5}
The word \emph{coarsening} can suggest compression. In the exact coordinate theorem it means regrouping complete packet vectors. Replacing each packet by one scalar sum is a different operation and really does lose information.

For $x=(1,2,3,4)$, the fine packets are the four coordinate-supported vectors. Regrouping gives $(1,0,3,0)$ and $(0,2,0,4)$. Every original coordinate remains available in one of these two vectors.

\vfig[.98]{regrouping_not_averaging}
\platecaption{The left and middle panels store the same four independent coordinates. The right panel instead applies a $2\times4$ numerical incidence matrix to four scalar entries. That is not the range-restricted exact coordinate map.}

Define
\[
X_N^{\rm pkt}=\bigoplus_j P_{N,j}X,
\qquad
\mathcal A_Nx=(P_{N,j}x)_j,
\qquad
\mathcal S_N(x_j)_j=\sum_jx_j.
\]
Completeness and mutual annihilation give both inverse laws. Consequently,
\[
J=\mathcal A_4\mathcal S_2,\qquad C=\mathcal A_2\mathcal S_4,
\qquad CJ=I,\quad JC=I.
\]
The apparently larger list of fine packets simply has smaller individual packet spaces.

\head{A kernel witnesses genuine information loss}
For the scalar incidence matrix $B$ drawn on the right,
\[
B(1,0,-1,0)^\mathsf T=0.
\]
The zero vector and $(1,0,-1,0)$ therefore have the same two scalar summaries. No inverse can recover the original four scalars from that output alone.

\takeaway{Ask what each packet stores. A list of full packet vectors can be exactly invertible. A list of packet means or totals generally is not.}

\vpage{5}{The operator changes coordinates too}{MODEL A / COHERENCE AND NATURALITY}{Chapter II, mathematical text: Theorem 3.5, Theorem 3.8}
Use a finite periodic carrier $X=\QQ^8$ for this example. Let $Se_j=e_{j+1\bmod8}$ and let
\[
T=I+2S+S^2.
\]
This is a concrete cyclic model, not an unannounced boundary convention for the finitely supported signal on Plate II.20. Write $R_N$ for the permutation that groups coordinates in increasing residue order modulo $N$.

\vfig[.99]{operator_naturality}
\platecaption{Actual $8\times8$ matrices. Dots are zero entries. The other entries are exact integers. Heavy lines mark the packet blocks: size four at level two and size two at level four.}

The reordered operators and change of coordinates are
\[
\widehat T_N=R_NTR_N^{-1},\qquad J_{4,2}=R_4R_2^{-1}.
\]
Multiplication cancels the intermediate coordinate maps:
\[
J_{4,2}\widehat T_2=\widehat T_4J_{4,2},\qquad
J_{8,4}J_{4,2}=J_{8,2}.
\]
Thus it makes no difference whether the operator is applied before or after a consistent coordinate change. One must transport the operator, not merely reuse the same numerical matrix under a new ordering.

\head{Refined blocks still reconstruct the same map}
A coarse block expands as
\[
P_bTP_a=\sum_{\beta\succ b,\,\alpha\succ a}P_\beta TP_\alpha.
\]
The terms are maps embedded in the original source and target spaces. Adding their norms, or adding their matrix entries after forgetting the embeddings, is not this operator identity.

\takeaway{The fine matrix contains a different block organization of the same operator. An independently guessed coarse matrix is a separate approximation problem, measured later by a defect.}

\vpage{6}{Selectors survive without division}{MODEL A / RESIDUE MASKS IN CHARACTERISTIC TWO}{Chapter II, mathematical text: Proposition 4.2}
Work first in $\mathbb F_2$, where $1+1=0$. The eight coefficients below are just a finite display of a coefficient sequence. The even and odd projectors are diagonal masks with entries zero and one.

\vfig[.95]{division_free_masks}
\platecaption{Each nonzero input coefficient is retained by exactly one mask. Reconstruction is exact over $\mathbb F_2$. No average and no root of unity is used.}

For a general coefficient ring and a positive integer $N$, the definition is
\[
(P_{N,j}x)_n=\begin{cases}x_n,&n\equiv j\pmod N,\\0,&\text{otherwise}.\end{cases}
\]
An address belongs to one residue class, so applying two masks either keeps that same coefficient or kills it. This proves
\[
P_{N,i}P_{N,j}=\delta_{ij}P_{N,j},\qquad \sum_jP_{N,j}=I.
\]
The argument is coordinatewise and does not require the coefficients to form a field.

\head{Why the usual even/odd average is a different formula}
Over a field of characteristic not two, parity extraction can be represented using $F(z)$ and $F(-z)$ and division by $2$. In characteristic two, $-z=z$ and $2=0$, so that formula cannot be used. The coefficient mask, however, still distinguishes even from odd \emph{exponents}.

This distinction also explains the later involution example. Failure to reduce a rational spectral projector containing $1/2$ does not make every integral coefficient-address selector disappear.

\takeaway{The basic object is the address selector. A Fourier average is one possible realization of it and has additional algebraic hypotheses.}

\vpage{7}{A Fourier certificate you can inspect}{MODEL A / AN EXACT NINE-ENTRY TABLE}{Chapter II, mathematical text: Theorem 4.3}
Take $N=3$ and work in $\mathbb F_7$. Set
\[
\omega=2,\qquad \omega^{-1}=4,\qquad \eta=3^{-1}=5.
\]
Here $2^3=1$ and neither $2$ nor $2^2$ is one. The unnormalized geometric sums are the entries in the first matrix.

\vfig[.97]{finite_fourier_table}
\platecaption{All entries are exact residues modulo seven. Multiplying the diagonal value $3$ by $5$ produces $1$. The lower calculation is a genuine failed Fourier certificate in a different coefficient ring.}

On the monomial $z^n$, define $UF(z)=F(2z)$. The proposed projector acts by
\[
5\sum_{k=0}^2 4^{jk}2^{nk}.
\]
The identity $2^3=1$ makes this depend only on $n\bmod3$. The nine-entry table therefore establishes the formula for every exponent once periodicity is used, rather than requiring a separate test at every coefficient.

\head{Order alone is not sufficient over every ring}
In $\mathbb F_3\times\mathbb F_3$, the element $(-1,1)$ has order two, but one component is a trivial character. Its nontrivial Fourier sum does not vanish in that component. Thus an element of order $N$ in a ring with zero divisors is not, by itself, the full orthogonality certificate.

\takeaway{The finite table is a checkable sufficient input. Over a field, a primitive root supplies it. Over a more general ring, inspect the actual table instead of inferring orthogonality from order alone.}

\vpage{8}{One address, two independent residues}{MODEL A / THE DYADIC AND ODD RECTANGLE}{Chapter II, mathematical text: Corollary 4.9}
Because $4$ and $3$ are coprime, a residue modulo $12$ is specified by a column $a=n\bmod4$ and a row $b=n\bmod3$. Each cell contains the unique combined address.

\vfig[.96]{crt_address_grid}
\platecaption{The twelve cells are a permutation of $0,\ldots,11$, not twelve extra degrees of freedom. The selected cell contains every integer congruent to $7$ modulo $12$.}

An explicit reconstruction formula is
\[
n\equiv9a+4b\pmod{12},
\]
because $9$ is one modulo four and zero modulo three, while $4$ has the opposite pair of residues. For $(a,b)=(3,1)$ this gives $31\equiv7$.

Applying the two masks in either order gives
\[
P_{12,7}=P_{4,3}P_{3,1}=P_{3,1}P_{4,3}.
\]
More generally, an address survives exactly when it satisfies both congruences. The coprime reconstruction identifies this intersection with one combined residue class.

\head{The hypothesis can be seen failing}
Replace $4$ and $3$ by $4$ and $6$. The requirements $n\equiv0\pmod4$ and $n\equiv1\pmod6$ are incompatible: the first makes $n$ even and the second makes it odd. A full rectangular address product no longer exists. Compatible pairs must agree modulo the greatest common divisor.

\takeaway{CRT separates address coordinates. It does not claim that dyadic and odd spectral projectors have the same denominator weights or the same lifting behavior.}

\vpage{9}{Which words can survive precision five?}{MODEL B / ORDERED TRANSFERS IN A FILTERED RING}{Chapter II, mathematical text: Theorem 5.10, Corollary 5.11}
Let $E_{ij}$ send the $j$th basis vector to the $i$th. In integral $2\times2$ matrices set
\[
a=2E_{01},\qquad b=4E_{10},\qquad A=a+b.
\]
Their certified dyadic weights are $1$ and $2$. At precision $L=5$, every word of weight at least five vanishes modulo $32$.

\vfig[.97]{weighted_word_cutoff}
\platecaption{All nonzero retained terms of the signed inverse sum, followed by the first discarded nonzero word. The empty word contributes $I$. Other words of weight below five can still be zero because $a^2=b^2=0$.}

The words are \emph{written matrix products}: in $ab$, the map $b$ acts first. Therefore
\[
ab=8E_{00},\qquad ba=8E_{11},\qquad aba=16E_{01}.
\]
The two products $ab$ and $ba$ cannot be identified merely because they have equal weight. Both their order and their source/target information matter.

The pruning rule supplies a sufficient reason for a term to vanish. It is not an assertion that every word below the threshold contributes nontrivially. Packet incompatibility and algebraic relations can remove additional words.

\takeaway{Length controls the sign. Cumulative weight controls precision. The ordered matrix product controls the actual contribution. None of these three can replace the others.}

\vpage{10}{An inverse computed in both orders}{MODEL B / A COMPLETE MODULO-32 CALCULATION}{Chapter II, mathematical text: Theorem 5.10, Corollary 5.11}
Keep the matrices from Plate II.9 and choose
\[
D=\begin{pmatrix}1&1\\0&1\end{pmatrix},\quad
D^{-1}=\begin{pmatrix}1&-1\\0&1\end{pmatrix},\quad
T=D(I+A)=\begin{pmatrix}5&3\\4&1\end{pmatrix}.
\]
The six retained contributions sum to
\[
S=I-a-b+ab+ba-aba
=\begin{pmatrix}9&-18\\-4&9\end{pmatrix}.
\]
The required inverse is $K=SD^{-1}$ in the quotient, not $D^{-1}S$.

\vfig[.98]{modular_inverse}
\platecaption{The displayed $K$ uses canonical residues from $0$ to $31$. Both left and right products are the identity modulo $32$. These are exact modular equalities, not floating-point residuals.}

One can verify the claim without trusting the word expansion:
\[
TK=\begin{pmatrix}129&64\\64&33\end{pmatrix}\equiv I,
\qquad
KT=\begin{pmatrix}65&32\\192&97\end{pmatrix}\equiv I\pmod{32}.
\]
The cancellation of the exterior factor $D$ explains the order. Since $T=D(I+A)$, multiply its inverse by $D^{-1}$ on the \emph{right} of the inverse of $I+A$.

\head{What the finite theorem does not require}
There is no claim here that the same integer series converges in the ordinary real matrix norm. The quotient kills all sufficiently deep products, so a finite sum suffices. The complete normed-ring theorem is a separate result with its own strict contraction hypothesis.

\takeaway{A correct noncommutative formula must work on both sides and preserve the order of its factors. Here the entire result can be checked by two integer matrix products.}

\vpage{11}{The residual tells you when to stop}{MODEL B / A CERTIFIED CUTOFF, NOT A HEURISTIC}{Chapter II, mathematical text: Theorem 5.10, Corollary 5.11, Section B.2}
For $S_N=\sum_{k=0}^{N-1}(-A)^k$, telescoping gives
\[
(I+A)S_N=S_N(I+A)=I-(-A)^N.
\]
Thus the exact residual is known. To certify precision $L$, it is enough to show that every residual entry is divisible by $2^L$.

\vfig[.94]{inverse_residual_depth}
\platecaption{The plotted values are the minimum dyadic valuations of the nonzero entries of $(-A)^N$. The dashed line is the requested cutoff $L=5$. This is an exact valuation profile, not a measured numerical convergence curve.}

In this particular model, $A^2=8I$. Therefore
\[
v_2(A^{2k})=3k,\qquad v_2(A^{2k+1})=3k+1,
\]
where the matrix valuation means the minimum entry valuation. Three geometric terms leave a residual of valuation four and do not suffice modulo $32$. Four terms leave valuation six and do suffice.

The generic positive-weight estimate is more conservative. If every edge weight is at least $m>0$, a word of length $N$ has weight at least $mN$. Hence the cutoff
\[
N=\left\lceil\frac Lm\right\rceil
\]
is sufficient. Here the certified minimum weight is $m=1$, so the generic choice is five. The extra identity $A^2=8I$ justifies stopping at four in this example.

\takeaway{A structural identity may improve a cutoff. An observed small real norm or an arbitrarily larger claimed minimum edge weight does not supply that certificate.}

\vpage{12}{Walks are not simple paths}{MODEL C / A GRAPH WITH PARALLEL EDGES}{Chapter II, mathematical text: Theorem 6.6}
Consider the directed graph below. The loop $\ell$ has weight one, the two distinct edges $a,b$ from $0$ to $1$ each have weight two, and the return edge $r$ has weight one.

\vfig[.96]{weighted_multigraph}
\platecaption{Every word on the right is read chronologically from left to right as an edge traversal. This is intentionally different from the right-to-left action of a written operator product in Plate II.9.}

A walk may revisit a vertex and reuse an edge. At total weight five there are six walks from $0$ to $1$. Two use three loops followed by either $a$ or $b$. Four travel to $1$, return to $0$, and travel to $1$ again, with two independent choices of parallel edge.

Their multiplicities enter the polynomial adjacency matrix:
\[
A(t)=\begin{pmatrix}t&2t^2\\t&0\end{pmatrix}.
\]
The coefficient $2$ remembers the two distinct edges. Replacing it by $1$ would count a different graph.

\head{The first-edge decomposition}
Let $a_n(s,t)$ count walks of total weight $n$ from $s$ to $t$. Removing the first edge gives
\[
a_n(s,t)=\mathbf1_{s=t,\,n=0}
 +\sum_{\substack{e:\,\mathrm{source}(e)=s\\w(e)\le n}}
 a_{n-w(e)}(\mathrm{target}(e),t).
\]
The empty walk is allowed precisely on the diagonal at weight zero. Positive weights ensure that every contributing walk has length at most $n$, so enumeration at each fixed weight is finite.

\takeaway{The graph series counts unsigned walks, including repetitions and parallel-edge multiplicity. It does not count simple paths, and it does not know whether an associated operator product vanishes or cancels.}

\vpage{13}{A finite graph produces its denominator}{MODEL C / AUTOMATIC RATIONALITY AND RECURRENCE}{Chapter II, mathematical text: Theorem 6.6, Corollary 6.7}
For the graph on Plate II.12, compute the determinant and adjugate rather than guessing a recurrence:
\[
I-A(t)=\begin{pmatrix}1-t&-2t^2\\-t&1\end{pmatrix},
\quad Q(t)=1-t-2t^3,
\quad \operatorname{adj}(I-A)_{0,1}=2t^2.
\]
Since $Q(0)=1$, its inverse exists as a formal power series with integer coefficients.

\vfig[.99]{zeta_recurrence}
\platecaption{The coefficients are independently computed by the first-edge recurrence and checked against the rational series through weight $24$. Literal edge-walk enumeration supplies a third check through weight $10$.}

From $QZ_{0,1}=2t^2$ one obtains
\[
a_0=a_1=0,\quad a_2=2,\qquad a_n=a_{n-1}+2a_{n-3}\quad(n\ge3).
\]
The numerator is responsible for the initial exception. Omitting it would incorrectly force the entire series to vanish.

\head{Three boundaries of this calculation}
This is a \emph{formal} series identity. No analytic convergence at a chosen real or complex value of $t$ is needed. The canonical determinant need not be the minimal denominator after cancellation. Finally, positive weights prevent infinitely many walks at a fixed weight: a single weight-zero loop would already give the infinitely many distinct walks $1,\ell,\ell^2,\ldots$ at weight zero.

\takeaway{Finite positive-weight graph data determine a rational series and an eventual recurrence. The proof does not begin by assuming a certificate $QZ=P$ that already contains that conclusion.}

\vpage{14}{Separate before dividing}{MODEL D / AN EVALUATED B\'EZOUT CERTIFICATE}{Chapter II, mathematical text: Theorem 7.5, Definition 7.6}
Use the integral involution
\[
A=\begin{pmatrix}1&1\\0&-1\end{pmatrix},\qquad A^2=I.
\]
With $F=A-I$ and $G=A+I$, the identities
\[
-F+G=2I,\qquad FG=0
\]
are an explicit evaluated B\'ezout certificate. Here all the required commutations follow because these operators are polynomials in the same matrix.

\vfig[.96]{bezout_normalization}
\platecaption{The denominator-free numerator is quasi-idempotent: $N^2=2N$. Division by the nonzero scalar $2$ gives an idempotent over the fraction field. It does not yet make the projector integral on the selected lattice.}

The numerator $N=G$ acts by $2$ on $\ker(A-I)$ and by zero on $\ker(A+I)$. Thus
\[
E_+=\frac12N
=\begin{pmatrix}1&1/2\\0&0\end{pmatrix}
\]
is the required separator. The complementary projector is $E_-=I-E_+$.

\head{A denominator bound need not be minimal}
The same projector can be presented with the inflated certificate $-2F+2G=4I$. Its numerator is $2G$ and its normalization is still $E_+$. This supplies a two-digit clearing bound at $2$, although one digit is enough. Thus the valuation of a supplied certificate gives an upper bound, not automatically the least possible weight.

For the original sign factors the concrete resultant is $\operatorname{Res}(x-1,x+1)=2$ and the bound is sharp. That equality is special to this displayed calculation, not a universal minimality theorem for every resultant certificate.

\takeaway{First check the evaluated identities and the numerator. Then normalize in the fraction field. Finally ask what denominator is needed on the particular lattice.}

\vpage{15}{Where the dyadic digit is lost}{MODEL D / THE EIGENLATTICE HAS INDEX TWO}{Chapter II, mathematical text: Proposition 9.3}
The eigenspaces are spanned by $(1,0)$ and $(-1,2)$. Their integral eigenvectors do not form an integral basis of $\ZZ^2$. The change-of-basis determinant is two.

\vfig[.94]{dyadic_eigenlattice}
\platecaption{An ordinary integer-lattice picture, not a drawing of the topology of $\QQ_2$. Filled dots have even second coordinate. Unfilled dots lie in the other coset. The hollow point at $(1/2,0)$ records the image of an integral vector under $E_+$.}

Write
\[
B=\begin{pmatrix}1&-1\\0&2\end{pmatrix},\qquad
AB=B\diag(1,-1),\qquad
E_+=B\diag(1,0)B^{-1}.
\]
The index-two eigenlattice explains the denominator in $B^{-1}$ and in the projector. Over the localized lattice $\Lambda_p=\ZZ_{(p)}^2$, the number $2$ is a unit for every odd prime $p$. Hence the local denominator weight is zero there and one at $p=2$.

This is a \emph{local} assertion. Odd localization does not change the fact that the displayed eigenvectors span an index-two sublattice inside the ordinary integer lattice.

\head{The special fiber does not retain the rational splitting}
Modulo two, $x-1=x+1$ and $A$ becomes a nontrivial unipotent matrix. The entries $1/2$ in $E_\pm$ cannot be reduced modulo two. By contrast, the coefficient selectors on Plate II.6 use only zero and one and remain available.

\takeaway{The obstruction is in integral separation, not in the existence of the two rational eigenspaces. The lattice is part of the question.}

\vpage{16}{One extra digit is genuinely necessary}{MODEL D / A SHARP PRECISION-DELAY EXAMPLE}{Chapter II, mathematical text: Theorem 7.7, Proposition 9.3}
The projector acts by
\[
E_+(x,y)=(x+y/2,0).
\]
Its factor $1/2$ can consume one dyadic digit. The claim concerns differences: if the inputs agree modulo $2^{m+1}\Lambda$, their projected differences lie in $2^m\Lambda$.

\vfig[.97]{precision_delay}
\platecaption{For $m=3$, inputs differing by $(0,8)$ agree modulo eight, but their projected outputs differ by $(4,0)$ and do not. Raising the input agreement to modulo sixteen makes the projected outputs agree modulo eight.}

The example works at every positive depth:
\[
E_+(0,2^m)=(2^{m-1},0)\notin2^m\Lambda,
\qquad
E_+(0,2^{m+1})=(2^m,0)\in2^m\Lambda.
\]
Thus a zero-delay promise fails, while one extra input digit is sufficient. This witnesses sharpness for the specified projector and lattice.

\head{The general calculation is just as short}
If $\pi^sE\Lambda\subseteq\Lambda$ and $x-y=\pi^{m+s}v$, then
\[
E(x-y)=\pi^m(\pi^sEv)\in\pi^m\Lambda.
\]
No rounding procedure is hiding in this identity. Nor does it authorize reducing $E$ modulo $\pi$ when its entries are nonintegral. Individual projected points may lie in a fractional lattice. The theorem controls their difference at the stated integral scale.

\head{Accuracy and branching are different views}
The later Smith example records how many inputs share one output at finite depth. Here we ask how accurately an input must be known to determine a projected output difference. Both involve directional divisibility, but they are different questions.

\takeaway{The phrase ``one lost digit'' has an exact meaning: both a sufficiency theorem and an explicit pair of inputs showing that a smaller delay fails.}

\vpage{17}{See every Smith fiber modulo eight}{MODEL E / TWO DIRECTIONAL WEIGHTS}{Chapter II, mathematical text: Theorem 10.6}
Consider the actual map of finite sets
\[
(\ZZ/8\ZZ)^2\longrightarrow(\ZZ/8\ZZ)^2,
\qquad (x,y)\longmapsto(2x,4y).
\]
The first coordinate of the output must be even. The second must be zero or four. Those are precisely the attainable targets in the right panel.

\vfig[.95]{smith_fibre_grid}
\platecaption{All sixty-four input points and all sixty-four possible target addresses are shown. The eight filled target points are the image. Every attainable target has exactly eight preimages. The other fifty-six targets have none.}

For a specific attainable target $(2,4)$,
\[
2x\equiv2\pmod8\iff x\in\{1,5\},\qquad
4y\equiv4\pmod8\iff y\in\{1,3,5,7\}.
\]
The coordinate choices are independent, giving $2\cdot4=8$ solutions. The entire carrier is partitioned consistently:
\[
8\text{ attainable targets}\ \times\ 8\text{ inputs per target}=64.
\]
For an inadmissible target, such as $(1,4)$ or $(2,2)$, the fiber is empty. The same nonzero cardinality formula must not be applied to it.

\head{The arithmetic behind the picture}
Multiplication by $p^s$ modulo $p^L$ has kernel size $p^{\min(s,L)}$. Its image consists exactly of residues whose canonical representatives are divisible by that same power. The displayed model has $p=2$, $L=3$ and weights $(1,2)$.

\takeaway{A directional weight has two consequences: a divisibility condition on the target and a conditional multiplicity of inputs. Solvability comes before counting.}

\vpage{18}{Transport the whole solution fiber}{MODEL E / EXPLICIT SMITH COORDINATES}{Chapter II, mathematical text: Theorem 10.6}
The preceding diagonal system is connected to a nondiagonal one by the unimodular change of variables
\[
M=\begin{pmatrix}2&2\\0&4\end{pmatrix},\quad
V=\begin{pmatrix}1&-1\\0&1\end{pmatrix},\quad
MV=\diag(2,4)=D.
\]
Since $\det V=1$, $x=Vy$ is bijective modulo eight. Thus $Mx=b$ is equivalent to $Dy=b$.

\vfig[.95]{smith_coordinate_transport}
\platecaption{The eight solutions for the target $(2,4)$, before and after the coordinate shear. Coordinates are residues in $\{0,\ldots,7\}$. No interpolation on this modular grid is intended.}

For example, the diagonal solution $y=(1,3)$ corresponds to
\[
x=Vy=(-2,3)\equiv(6,3)\pmod8,
\]
and $M(6,3)^\mathsf T=(18,12)^\mathsf T\equiv(2,4)^\mathsf T$. Applying the inverse shear recovers the original solution.

\head{Two boundary cases worth keeping visible}
A unit diagonal coordinate has exactly one solution for every target. A literal zero coordinate has all eight possible inputs when its target is zero, and no input when the target is nonzero. A zero coordinate is not inverted or silently omitted.

The displayed diagonalization is verified by multiplication. The theorem takes diagonal data as input. A Smith normal form for an arbitrary matrix or general DVR requires a separate structure theorem.

\takeaway{An invertible change of coordinates transports existence, every solution, and the fiber count. The target must be transformed as prescribed when checking admissibility.}

\vpage{19}{Precision exposes directional freedom}{MODEL E / A DEPTH-BY-DEPTH FIBER BUDGET}{Chapter II, mathematical text: Theorem 10.6}
Take $p=2$ and the supplied diagonal weights $(s_0,s_1,s_2)=(0,1,3)$. At depth $L$, an admissible target has
\[
2^{W(L)}\text{ solutions},\qquad
W(L)=\min(0,L)+\min(1,L)+\min(3,L).
\]
Each direction contributes only the amount visible at the current depth.

\vfig[.93]{smith_depth_budget}
\platecaption{The stacked height records the exponent $W(L)$. The label above each bar records the actual number of solutions. The zero-weight direction contributes no branching. This is a finite cardinality calculation, not an information-theoretic assumption about random inputs.}

The values are
\[
\begin{array}{c|rrrrrr}
L&1&2&3&4&5&6\\\hline
W(L)&2&3&4&4&4&4\\
\#\text{ solutions}&4&8&16&16&16&16.
\end{array}
\]
Once $L\ge3$, the nonzero weights have been fully exposed and the admissible fiber cardinality stabilizes at sixteen. The input space still grows as $2^{3L}$. The number of attainable targets grows accordingly.

At depth six, $(1,2,8)$ is admissible and gives sixteen solutions. Replacing the last coordinate by four makes the target inadmissible and gives zero solutions, even though the weights and depth are unchanged.

\head{What changes with a zero diagonal factor?}
A literal zero factor contributes $L$, not a fixed eventual weight, and requires a zero target coordinate. Its fiber therefore keeps growing with the modulus. That case must not be mistaken for a finite nonzero Smith weight.

\takeaway{The scalar sum of weights gives the count, but the separate weights give the admissibility conditions. Keeping only a determinant valuation can hide which target coordinate fails.}

\vpage{20}{Filter first or split first: same output}{MODEL F / A FINITE POLYPHASE REALIZATION}{Chapter II, mathematical text: Section B.1}
Use the finitely supported signal $x=(1,2,3,4,5,6)$ and the filter $h=(1,2,1)$. Outside the displayed support, the input is zero. Let $e_m=x_{2m}$ and $o_m=x_{2m+1}$.

\vfig[.96]{finite_polyphase_filter}
\platecaption{Exact integer convolution, including both output tail coefficients. The two bottom rows are the even and odd output phases. Interleaving them reconstructs the full eight-term output.}

If $w$ denotes one delay on a phase sequence, the filter becomes
\[
\begin{pmatrix}y_{\rm even}\\y_{\rm odd}\end{pmatrix}
=
\underbrace{\begin{pmatrix}1+w&2w\\2&1+w\end{pmatrix}}_{H_{\rm phase}(w)}
\begin{pmatrix}e\\o\end{pmatrix}.
\]
The cross-phase delay is not optional. A unit shift sends an odd sample to the next even position, while it sends an even sample to the odd position in the same phase block.

Here $e=1+3w+5w^2$ and $o=2+4w+6w^2$, so direct polynomial multiplication gives
\[
y_{\rm even}=1+8w+16w^2+17w^3,
\qquad
y_{\rm odd}=4+12w+20w^2+6w^3.
\]
This is the same result as ordinary convolution followed by residue extraction. No numerical transform or limiting argument is needed because every sequence is finitely supported.

\takeaway{A phase matrix is not just a sparsity pattern. Its delay factors carry the quotient-and-remainder information needed when a shift crosses a phase boundary.}

\vpage{21}{Two Smith constructions, different roles}{MODEL G / A NONLINEAR PHASE ORBIT}{Chapter II, mathematical text: Section 12.3, Section 12.3, Smith-phase identity (p. 43)}
Smith \emph{normal form} concerns linear congruences. The Smith \emph{impedance chart} uses the nonlinear coordinate change
\[
\Gamma=\frac{z-1}{z+1},\qquad z=\frac{1+\Gamma}{1-\Gamma}.
\]
They share a name, not a mathematical definition. Sample the reflection circle at the exact algebraic points $\Gamma_k=\tfrac12 e^{ik\pi/4}$.

\vfig[.96]{smith_chart_orbit}
\platecaption{Matching labels identify the same point in the two coordinates. All eight points are computed algebraically. Only their placement on the page is rendered numerically. The outer left circle is $|\Gamma|=1$, not part of the sampled radius-$1/2$ orbit.}

The half-turn $\Gamma\mapsto-\Gamma$ gives
\[
z\longmapsto\frac1z,\qquad z\ne0,-1.
\]
Thus the labels $0$ and $4$ are $3$ and $1/3$. Label $2$ is $(3+4i)/5$, and the opposite point is its reciprocal. The domain restrictions come from the forward and inverse coordinate denominators.

These sampled phase points are not linear idempotent ranges. Refining from two to four to eight samples inserts more points on the same orbit. This does not imply a linear packet decomposition of impedance values.

\takeaway{The chart is a geometric realization of phase action. Its nonlinear coordinate sensitivity near $\Gamma=\pm1$ is not the Smith congruence-fiber count and not a new linear projector theorem.}

\vpage{22}{A small defect can grow under powers}{MODEL H / EXACT TELESCOPING, EXPLICIT GROWTH}{Chapter II, mathematical text: Theorem 11.1, Corollary 11.3}
Choose an invertible coordinate shear but independently keep the same numerical operator at both levels:
\[
C=\begin{pmatrix}1&\varepsilon\\0&1\end{pmatrix},\quad
T_f=T_c=\diag(1,2),\quad \varepsilon=\frac1{16}.
\]
Then $\Delta=CT_f-T_cC=\varepsilon E_{01}$. Exact coordinate invertibility has not forced these independently selected operators to intertwine.

\vfig[.95]{power_defect_growth}
\platecaption{Actual defects and certified upper bounds in the induced infinity norm, both divided by $\varepsilon$. All heights are exact integers. The upper estimate is valid but need not be tight.}

Direct multiplication gives
\[
CT_f^n-T_c^nC=\varepsilon(2^n-1)E_{01}.
\]
The general telescoping formula explains every contribution:
\[
CT_f^n-T_c^nC
=\sum_{k=0}^{n-1}T_c^{n-1-k}\Delta T_f^k.
\]
Since both operator norms are two, the source theorem bounds the error by $n2^{n-1}\varepsilon$. It does not promise a uniform-in-$n$ small error merely because $\varepsilon$ is small.

Had we chosen the transported operator $T_c'=CT_fC^{-1}$, the defect would be zero and every polynomial action would intertwine exactly. That is a different choice of coarse operator, not a stronger conclusion from the faulty one.

\takeaway{Separate coordinate coherence from operator consistency. A defect from one step is propagated by the dynamics on both sides. The telescoping identity tells you exactly how.}

\vpage{23}{Resolvents amplify errors near poles}{MODEL H / A DOMAIN-SENSITIVE COMPARISON}{Chapter II, mathematical text: Theorem 11.5}
For the same model, take $z\notin\{1,2\}$ and form the actual inverse matrices
\[
R_f=(zI-T_f)^{-1},\qquad R_c=(zI-T_c)^{-1}.
\]
Their off-diagonal intertwining error is
\[
CR_f-R_cC=R_c\Delta R_f
=\frac{\varepsilon}{(z-1)(z-2)}E_{01}.
\]
Both inverse hypotheses matter. Neither resolvent is asserted at a pole.

\vfig[.94]{resolvent_amplification}
\platecaption{The bars show exact amplification factors for five rational choices $z>2$. Rows are separate samples, not equally spaced positions on the $z$ axis. The coefficient $\varepsilon=1/16$ has been divided out.}

At $z=3$ the amplification factor is only $1/2$. At $z=201/100$, it is $10000/101$, so a small one-step defect can create a much larger resolvent defect. This follows from the exact formula rather than from unstable floating-point subtraction of two nearly singular inverses.

The certified general norm estimate is
\[
\|CR_f-R_cC\|
\le \|R_c\|\,\|\Delta\|\,\|R_f\|.
\]
The paper also proves a one-sided-witness version: a left inverse for the coarse shift and a right inverse for the fine shift suffice for the algebraic identity. That does not turn arbitrary one-sided witnesses into a claim of two-sided resolvent existence.

\takeaway{A small local defect is meaningful only together with control of the surrounding operators or inverses. Near a pole, the domain and amplification factors are part of the theorem, not technical decoration.}

\vpage{24}{Read the certificates with the pictures}{PROOF BOUNDARY / HOW TO USE THIS GUIDE}{Formal evidence: chapter-wide correspondence and the 25-row foundational contract}
The foundational formal contract has twenty-five items: twenty-one numbered main-text definitions and results, one Smith-phase identity, and three algorithm specifications. The source archive preserves its evidence. The broader constructions have a separate correspondence and replay record. The table below describes the foundational interfaces; the chapter-wide map identifies the broader constructions. The models illustrate the mathematics, and their finite checks have recorded ranges.

\input{generated/current_visual_scope.tex}

{\small
\begin{tabularx}{\linewidth}{@{}p{.28\linewidth}YY@{}}\toprule
\textbf{Inherited interface} & \textbf{Inputs and checked scope} & \textbf{Boundary of that interface}\\\midrule
Exact refinement & Complete projector families and a parent map when formulas relating parents to children are used. & A canonical refinement for every arbitrary pair of decompositions.\\
Fourier realization & A finite orthogonality table, or field/root hypotheses that provide it. & Orthogonality from multiplicative order alone over every ring.\\
Weighted inversion & A two-sided inverse of $D$, a multiplicative filtration, and positive certified gains. & Convergence in every possible complete filtered algebra.\\
Graph rationality & A finite multigraph with positive weights, together with its actual determinant and adjugate. & Simple-path counts, analytic convergence everywhere, or operator noncancellation.\\
Spectral separation & Evaluated B\'ezout identities and explicit lattice-preservation data. & Automatic construction or minimality of every resultant certificate.\\
Smith congruences & Concrete diagonal data modulo $p^L$, with $p>1$. & An automatically constructed general-DVR Smith normal form.\\
Linear defects & Actual typed maps and the stated inverse or norm witnesses. & Small propagated errors independent of time or distance from a pole.\\\bottomrule
\end{tabularx}}

\head{Reproduce the explanatory calculations}
The preserved model source contains \textbf{118 named exact checks}. Each finite range is recorded in the file \texttt{EXACT\_CHECKS.json}. They cover masks, CRT labels, modular inverse products, graph enumeration, rational recurrences, lattice denominators, complete finite fibers, phase identities, and defect formulas. Graphics are generated from those same exact model definitions. Numerical placement of geometric points is only a rendering step.

The computational report records 48 passing Python tests, 118 exact visual checks, 22 byte-matching regenerated vector plates, 13 regenerated PDF/PNG figure pairs, and nine matching numerical summary rows. Each record retains its tested range and source identity. Lean replay and source coverage are documented separately.

\takeaway{Use the guide to understand what a theorem says and why its hypotheses matter. Use the theorem map and preserved proof sources to identify the precise formal result. Neither a picture nor a finite test expands that result's scope.}

\endgroup
\clearpage
\bibliographystyle{plain}\bibliography{references}
\end{document}
